\PassOptionsToPackage{unicode}{hyperref}
\PassOptionsToPackage{hyphens}{url}
\PassOptionsToPackage{dvipsnames,svgnames,x11names}{xcolor}
\documentclass[10pt,letterpaper]{article}

\usepackage{xcolor}
\usepackage[margin=2.0cm]{geometry}
\usepackage{lmodern}
\usepackage{amsmath,amssymb,amsfonts,mathtools}
\usepackage{textcomp}
\usepackage{microtype}
\usepackage{parskip}
\usepackage{longtable,booktabs,array,calc}
\usepackage{caption}
\captionsetup[table]{skip=6pt,font=small,labelfont=bf}
\captionsetup[figure]{font=small,labelfont=bf}
\newcounter{none}
\usepackage{etoolbox}
\usepackage{graphicx}
\usepackage{fancyvrb}
\newcommand{\VerbBar}{|}
\newcommand{\VERB}{\Verb[commandchars=\\\{\}]}
\DefineVerbatimEnvironment{Highlighting}{Verbatim}{commandchars=\\\{\},fontsize=\footnotesize}
\RecustomVerbatimEnvironment{verbatim}{Verbatim}{fontsize=\footnotesize}
\newenvironment{Shaded}{}{}
\newcommand{\AlertTok}[1]{\textcolor[rgb]{0.80,0.00,0.00}{\textbf{#1}}}
\newcommand{\AnnotationTok}[1]{\textcolor[rgb]{0.38,0.63,0.69}{\textbf{\textit{#1}}}}
\newcommand{\AttributeTok}[1]{\textcolor[rgb]{0.49,0.56,0.16}{#1}}
\newcommand{\BaseNTok}[1]{\textcolor[rgb]{0.25,0.63,0.44}{#1}}
\newcommand{\BuiltInTok}[1]{\textcolor[rgb]{0.00,0.50,0.00}{#1}}
\newcommand{\CharTok}[1]{\textcolor[rgb]{0.25,0.44,0.63}{#1}}
\newcommand{\CommentTok}[1]{\textcolor[rgb]{0.38,0.63,0.69}{\textit{#1}}}
\newcommand{\CommentVarTok}[1]{\textcolor[rgb]{0.38,0.63,0.69}{\textbf{\textit{#1}}}}
\newcommand{\ConstantTok}[1]{\textcolor[rgb]{0.53,0.00,0.00}{#1}}
\newcommand{\ControlFlowTok}[1]{\textcolor[rgb]{0.00,0.44,0.13}{\textbf{#1}}}
\newcommand{\DataTypeTok}[1]{\textcolor[rgb]{0.56,0.13,0.00}{#1}}
\newcommand{\DecValTok}[1]{\textcolor[rgb]{0.25,0.63,0.44}{#1}}
\newcommand{\DocumentationTok}[1]{\textcolor[rgb]{0.73,0.13,0.13}{\textit{#1}}}
\newcommand{\ErrorTok}[1]{\textcolor[rgb]{0.80,0.00,0.00}{\textbf{#1}}}
\newcommand{\ExtensionTok}[1]{#1}
\newcommand{\FloatTok}[1]{\textcolor[rgb]{0.25,0.63,0.44}{#1}}
\newcommand{\FunctionTok}[1]{\textcolor[rgb]{0.02,0.16,0.49}{#1}}
\newcommand{\ImportTok}[1]{\textcolor[rgb]{0.00,0.50,0.00}{\textbf{#1}}}
\newcommand{\InformationTok}[1]{\textcolor[rgb]{0.38,0.63,0.69}{\textbf{\textit{#1}}}}
\newcommand{\KeywordTok}[1]{\textcolor[rgb]{0.00,0.44,0.13}{\textbf{#1}}}
\newcommand{\NormalTok}[1]{#1}
\newcommand{\OperatorTok}[1]{\textcolor[rgb]{0.40,0.40,0.40}{#1}}
\newcommand{\OtherTok}[1]{\textcolor[rgb]{0.00,0.44,0.13}{#1}}
\newcommand{\PreprocessorTok}[1]{\textcolor[rgb]{0.74,0.48,0.00}{#1}}
\newcommand{\RegionMarkerTok}[1]{#1}
\newcommand{\SpecialCharTok}[1]{\textcolor[rgb]{0.25,0.44,0.63}{#1}}
\newcommand{\SpecialStringTok}[1]{\textcolor[rgb]{0.73,0.40,0.53}{#1}}
\newcommand{\StringTok}[1]{\textcolor[rgb]{0.25,0.44,0.63}{#1}}
\newcommand{\VariableTok}[1]{\textcolor[rgb]{0.10,0.09,0.49}{#1}}
\newcommand{\VerbatimStringTok}[1]{\textcolor[rgb]{0.25,0.44,0.63}{#1}}
\newcommand{\WarningTok}[1]{\textcolor[rgb]{0.38,0.63,0.69}{\textbf{\textit{#1}}}}

\IfFileExists{xeCJK.sty}{%
  \usepackage{xeCJK}
  \IfFontExistsTF{Songti SC}{%
    \setCJKmainfont{Songti SC}
    \setCJKmonofont{Songti SC}
  }{}
  \IfFontExistsTF{Apple SD Gothic Neo}{%
    \xeCJKDeclareSubCJKBlock{Hangul}{"AC00 -> "D7AF}
    \setCJKmainfont[Hangul]{Apple SD Gothic Neo}
  }{}
}{}

\usepackage{newunicodechar}
\newunicodechar{µ}{\ifmmode\mu\else\textmu\fi}
\newunicodechar{°}{\ifmmode^\circ\else\textdegree\fi}
\newunicodechar{·}{\ifmmode\cdot\else\textperiodcentered\fi}
\newunicodechar{²}{\ensuremath{^{2}}}
\newunicodechar{³}{\ensuremath{^{3}}}
\newunicodechar{¹}{\ensuremath{^{1}}}
\newunicodechar{⁴}{\ensuremath{^{4}}}
\newunicodechar{⁶}{\ensuremath{^{6}}}
\newunicodechar{⁻}{\ensuremath{^{-}}}
\newunicodechar{₀}{\ensuremath{_{0}}}
\newunicodechar{ₙ}{\ensuremath{_{n}}}
\newunicodechar{Δ}{\ensuremath{\Delta}}
\newunicodechar{Ω}{\ensuremath{\Omega}}
\newunicodechar{α}{\ensuremath{\alpha}}
\newunicodechar{β}{\ensuremath{\beta}}
\newunicodechar{δ}{\ensuremath{\delta}}
\newunicodechar{ε}{\ensuremath{\varepsilon}}
\newunicodechar{θ}{\ensuremath{\theta}}
\newunicodechar{κ}{\ensuremath{\kappa}}
\newunicodechar{λ}{\ensuremath{\lambda}}
\newunicodechar{π}{\ensuremath{\pi}}
\newunicodechar{ρ}{\ensuremath{\rho}}
\newunicodechar{σ}{\ensuremath{\sigma}}
\newunicodechar{τ}{\ensuremath{\tau}}
\newunicodechar{→}{\ensuremath{\rightarrow}}
\newunicodechar{∎}{\ensuremath{\blacksquare}}
\newunicodechar{−}{\ensuremath{-}}
\newunicodechar{∧}{\ensuremath{\wedge}}
\newunicodechar{∫}{\ensuremath{\int}}
\newunicodechar{≈}{\ensuremath{\approx}}
\newunicodechar{≤}{\ensuremath{\leq}}
\newunicodechar{≥}{\ensuremath{\geq}}
\newunicodechar{⊆}{\ensuremath{\subseteq}}
\newunicodechar{✓}{\ensuremath{\checkmark}}
\newunicodechar{≠}{\ensuremath{\neq}}
\newunicodechar{↔}{\ensuremath{\leftrightarrow}}
\newunicodechar{≪}{\ensuremath{\ll}}
\newunicodechar{↛}{\ensuremath{\nrightarrow}}
\newunicodechar{⇒}{\ensuremath{\Rightarrow}}
\newunicodechar{∪}{\ensuremath{\cup}}
\newunicodechar{⟨}{\ensuremath{\langle}}
\newunicodechar{⟩}{\ensuremath{\rangle}}
\newunicodechar{・}{\textperiodcentered}

\usepackage{xurl}
\usepackage{bookmark}
\usepackage{hyperref}

\setcounter{secnumdepth}{-\maxdimen}
\providecommand{\tightlist}{\setlength{\itemsep}{1pt}\setlength{\parskip}{0pt}}

\definecolor{openieAmber}{HTML}{D97706}
\definecolor{openieAmberBg}{HTML}{FFFBEB}
\definecolor{openieAmberBorder}{HTML}{F59E0B}
\definecolor{openieSlate}{HTML}{334155}
\definecolor{openieDark}{HTML}{0F172A}
\definecolor{openieCodeBg}{HTML}{F8FAFC}
\definecolor{openieCodeBorder}{HTML}{CBD5E1}
\definecolor{openieBlue}{HTML}{0057B8}

\hypersetup{
  pdftitle={The Gap Between Possible and Deployed: A Cadence Model of Computer Intelligence and the Forcing Function of Automation Tools},
  pdfauthor={David Charlot, Open Interface Engineering},
  colorlinks=true,
  linkcolor=openieBlue,
  citecolor=openieBlue,
  urlcolor=openieBlue,
  pdfcreator={XeTeX / Tectonic — Open Interface Engineering}
}

\urlstyle{same}
\setlength{\emergencystretch}{3em}
\hfuzz=6pt
\sloppy

\newcommand{\openiefig}[3]{%
  \IfFileExists{#1}{%
    \includegraphics[width=#2\linewidth,height=0.70\textheight,keepaspectratio]{#1}%
  }{%
    \fcolorbox{openieCodeBorder}{openieCodeBg}{%
      \parbox{0.92\linewidth}{\footnotesize\ttfamily Figure asset embedded in archival PDF; interactive/vector source: \url{#3}}%
    }%
  }%
}

\title{\vspace{-1.2em}\textbf{The Gap Between Possible and Deployed: A Cadence Model of Computer Intelligence and the Forcing Function of Automation Tools}}
\author{\textbf{David Charlot, Open Interface Engineering}\\[0.15em]
{\small Research study $\cdot$ Open Interface Engineering $\cdot$ \href{https://research.openie.dev}{research.openie.dev}}}
\date{\small Revision \texttt{395f081} (4 November 2026) $\cdot$ Web: \href{https://research.openie.dev/papers/latency/}{research.openie.dev/papers/latency/} $\cdot$ Archival LaTeX: \href{https://research.openie.dev/pdfs/latex/latency.pdf}{/pdfs/latex/latency.pdf} (\href{https://research.openie.dev/pdfs/latex/latency.tex}{.tex})}

\begin{document}
\maketitle
\vspace{-0.6em}

% OpenIE Measurement Honesty Callout Box
\noindent\fcolorbox{openieAmberBorder}{openieAmberBg}{%
  \parbox{\dimexpr\linewidth-2\fboxsep-2\fboxrule\relax}{%
    \small\color{openieDark}%
    \textbf{Measurement.} Research study, not a journal final. Publication-to-deployment lags, quarterly cadence tests, corporate capital expenditures, remaining performance obligations, circular financing edges, grid power queues, and inference price series are compiled from dated \texttt{SEC} filings (\texttt{10-K}/\texttt{10-Q}), regulatory disclosures, and primary technical reports. Task-level dollar and joule figures in the cost model are analytical accounting estimates (\texttt{est\_j}) parameterized from published hardware specifications and literature measurements; no new hardware rail or wall meter reading is claimed (\texttt{measured\_j = None}, \texttt{board\_metered = false}). No fabricated citations.
  }%
}

\vspace{6pt}
% OpenIE Amber Left-Rule Summary Quote
\noindent{\color{openieAmber}\vrule width 3pt}\hspace{8pt}%
\parbox{\dimexpr\linewidth-11pt\relax}{%
  \small\itshape\color{openieSlate}%
  Methods that lower the cost of computer intelligence appear in public years before most products use them. This article measures that gap from dated sources, models how reporting and funding calendars set part of it, tests how much of the infrastructure spending is a closed loop, models the cost of a task, and states what data would refute the model.
}
\vspace{6pt}

{
\small
\setcounter{tocdepth}{2}
\tableofcontents
}
\vspace{4pt}
\hrule
\vspace{8pt}

\emph{18 October 2026. Sources are dated and sorted by evidence class in
the References. The propositions are derived by hand and have not been
checked by a proof assistant, and none of the numbers in Section 7 is
data. Quantities marked model output are computed from cited figures and
are not observations.}

\section{Abstract}\label{abstract}

New methods for running computer intelligence at lower cost are
published openly and then reach most products after a delay. This
article defines that delay as a measurable quantity, splits it into
technical, calendar and incentive parts, and models how a tool that
lowers the cost of automation changes what a seller can profitably
charge for. It then places the delay inside the present build-out of
infrastructure. It sets out what public filings document about money
moving between labs, clouds and chip makers, states a test that
separates a closed loop from outside demand, builds a cost model in
dollars and joules per task, compares open software methods and
alternative hardware with the incumbent stack, and maps where
capability, compute and power sit. The propositions are derived by hand,
and each prediction names the data that would refute it.

\textbf{In plain words}

\begin{enumerate}
\def\labelenumi{\arabic{enumi}.}
\tightlist
\item
  Open methods that make computer intelligence cheaper reach the same
  results on older and cheaper machines, and the wait between a
  published method and a shipped product is shorter for some methods
  than for others.
\item
  Most of what a task costs is the price of the machines and not the
  electricity, and prices paid by users have fallen far faster than
  hardware costs. The money behind the machines includes loops among a
  few labs, clouds and chip makers that filings document, and it also
  includes cash from other businesses and from customers outside those
  loops. Section 10 gives a test that separates the two with data.
\item
  Part of the wait comes from calendars, because companies decide in
  quarterly and yearly steps. When a tool makes the work cheap enough,
  waiting stops paying for the seller, and what is left to charge for is
  the cost of switching away and not the cost of computing.
\end{enumerate}

\section{Notation}\label{notation}

{\def\LTcaptype{none} % do not increment counter
{\footnotesize\setlength{\tabcolsep}{3.8pt}\renewcommand{\arraystretch}{1.15}
\begin{longtable}[]{@{}
  >{\raggedright\arraybackslash}p{(\linewidth - 2\tabcolsep) * \real{0.3125}}
  >{\raggedright\arraybackslash}p{(\linewidth - 2\tabcolsep) * \real{0.6875}}@{}}
\toprule\noalign{}
\begin{minipage}[b]{\linewidth}\raggedright
\textbf{Term}
\end{minipage} & \begin{minipage}[b]{\linewidth}\raggedright
\textbf{Meaning in this article}
\end{minipage} \\
\midrule\noalign{}
\endhead
\bottomrule\noalign{}
\endlastfoot
Task & One unit of useful work with a stated quality threshold \(q\),
for example one correct answer or one checked estimate. \\
Joules per task, \(e\) & Measured energy in joules to complete one task
at quality at least \(q\). The unit of account for computer intelligence
here. \\
Cost per task, \(c\) & Money cost of one task: \(c = \kappa e + h\),
where \(\kappa\) is the price of a joule and \(h\) is the hardware and
operations cost that does not fall with \(e\). \\
Capability frontier, \(C_f(t)\) & The lowest cost per task reachable by
any method that is publicly documented by time \(t\). \\
Deployed frontier, \(D(t)\) & The lowest cost per task offered in a
generally available product by time \(t\). \\
Latency, \(L\) & At a cost level \(c\), the time between the capability
frontier reaching \(c\) and the deployed frontier reaching \(c\).
Defined in Section 0. \\
Reporting period, \(T\) & The spacing of decision dates in a firm that
reports earnings, about a quarter. \\
Switching friction, \(s\) & The per-task cost a buyer bears to move to
another seller: integration, data transfer, retraining, contract terms,
default settings. \\
Adoption cost, \(K\) & One-time cost to rebuild an offering around a new
method. \\
Accelerator time per task, \(\tau\) & Utilized accelerator-hours that
one task occupies at quality at least \(q\). Section 12. \\
Capital-loaded price of a joule, \(\kappa^{*}\) & Electricity price plus
capital per utilized hour divided by facility power. Section 12. \\
Closed-loop share bound, \(\bar\varphi_s\) & Upper bound on the share of
a supplier's revenue that its own supplied capital could have funded.
Section 10. \\
Tolerated shortfall, \(w\) & The largest fall in reported earnings in
one period that a firm will accept, in the same units as \(K\). \\
\end{longtable}
}
}

\section{0. Definitions and scope}\label{definitions-and-scope}

Let \(c_C(t)\) be the lowest cost per task reachable at time \(t\) by
any publicly documented method at quality \(q\). It is non-increasing,
because documentation is not retracted. Let \(c_D(t)\) be the lowest
cost per task offered in a generally available product at time \(t\). It
is non-increasing for the same reason when a seller does not withdraw an
offer, and \(c_D(t) \ge c_C(t)\) by construction. For a cost level
\(c\), write \(t_C(c) = \inf\{t : c_C(t) \le c\}\) and
\(t_D(c) = \inf\{t : c_D(t) \le c\}\).

\textbf{Definition 1 (latency).} \(L(c) = t_D(c) - t_C(c) \ge 0\). The
question of the article is how large \(L\) is, what sets it, and what
moves it.

\textbf{Definition 2 (technical, calendar and incentive parts).} For a
firm \(i\) the latency is decomposed as
\(L_i = d + w^{cad}_i + w^{gate}_i\). Here \(d\) is the technical
readiness delay: the time engineers need to port, test and harden the
method. The term \(w^{cad}_i\) is calendar wait: time between readiness
and the next date on which the firm makes release decisions. The term
\(w^{gate}_i\) is incentive delay: time during which release has
positive long-run value but a reporting or funding rule blocks it. The
\textbf{artificial latency} is \(L_i - d\). The word is used in a narrow
sense: it is the part not explained by engineering time, and it is not a
claim about anyone's motive.

\textbf{Scope.} The article covers methods that change the cost of
computer intelligence: model architectures, number formats, kernels,
compilers, statistical inference procedures and energy measurement. It
also covers the infrastructure that runs them: the flows of capital and
compute commitments among labs, clouds and chip makers, the cost of an
accelerator-hour, open methods and alternative hardware that change the
accelerator time per task, and where capability, compute and power are
located. It does not forecast any company, price or share. Energy per
task in joules is the metric where a joule figure exists. Where a source
reports only dollars or seconds, the dollar or second figure is used and
labeled as such.

\section{1. Kinds of evidence}\label{kinds-of-evidence}

Every source in this article is sorted into one of three classes, and
the class is printed in the References.

\begin{itemize}
\tightlist
\item
  \textbf{Established by proof in the source.} The source contains a
  theorem or analytic derivation with its proof. A proof inside a paper
  is the author's proof, not a machine-checked one, and the article does
  not call it more. Court records and statistical releases are sorted as
  demonstrated, because they report dated facts and not derivations.
\item
  \textbf{Demonstrated.} The source reports a measurement, survey,
  experiment or dated release with a stated method. The result holds for
  that setting and may not carry elsewhere.
\item
  \textbf{Opinion.} The source is a forecast, a vendor statement without
  data, a practitioner narrative, or a review that summarizes other
  work. Reviews are sorted here because the article cannot tell from
  them what was measured.
\end{itemize}

Figures are quoted as the source states them, with the source's own
setting and caveats. Where a quantity is derived from cited figures, the
text calls it model output and shows the arithmetic or names the script
that computes it. Announcements of plans, targets and intended
investments are sorted as opinion, because they state intent and not
outcome, and filings with the regulator are sorted as demonstrated.
Every source was retrieved as full text and each cited figure was
located in that text. Sources in German, French, Spanish, Portuguese,
Italian, Korean, Japanese and Chinese were read in the original
language, and their figures are restated in English without change.

\section{2. What became possible, 2024 to
2026}\label{what-became-possible-2024-to-2026}

The list below is organized by the lever each method pulls on cost per
task. Each row gives the date of earliest public posting, the claim as
the authors report it, and the class.

\textbf{Table 1. Methods that change what a task costs.}

{\def\LTcaptype{none} % do not increment counter
{\footnotesize\setlength{\tabcolsep}{3.8pt}\renewcommand{\arraystretch}{1.15}
\begin{longtable}[]{@{}
  >{\raggedright\arraybackslash}p{(\linewidth - 6\tabcolsep) * \real{0.1667}}
  >{\raggedright\arraybackslash}p{(\linewidth - 6\tabcolsep) * \real{0.4815}}
  >{\raggedright\arraybackslash}p{(\linewidth - 6\tabcolsep) * \real{0.1204}}
  >{\raggedright\arraybackslash}p{(\linewidth - 6\tabcolsep) * \real{0.2315}}@{}}
\toprule\noalign{}
\begin{minipage}[b]{\linewidth}\raggedright
\textbf{Lever}
\end{minipage} & \begin{minipage}[b]{\linewidth}\raggedright
\textbf{Method and claim}
\end{minipage} & \begin{minipage}[b]{\linewidth}\raggedright
\textbf{Posted}
\end{minipage} & \begin{minipage}[b]{\linewidth}\raggedright
\textbf{Class}
\end{minipage} \\
\midrule\noalign{}
\endhead
\bottomrule\noalign{}
\endlastfoot
Architecture & Selective state space models and attention variants are
dual through structured semiseparable matrices, and the resulting layer
is reported 2 to 8 times faster than the earlier scan {[}38{]}. & 31 May
2024 & established by proof in the source \\
Architecture & Hybrid state space and attention models were released
with open weights by several groups across 2024 and 2025 {[}39{]}. &
2024 to 2025 & demonstrated \\
Number format & Ternary weights match 16-bit perplexity and end-task
scores from about 3 billion parameters at lower latency, memory and
energy, in the authors' experiments {[}40{]}. & 27 Feb 2024 &
demonstrated \\
Number format & Stable 4-bit pretraining of a 12 billion parameter
hybrid model on 10 trillion tokens, comparable to 8-bit results,
reported by the hardware vendor {[}41{]}. & Sep 2025 & demonstrated \\
Kernels & An attention kernel reaches up to 75\% of peak 16-bit
throughput on one accelerator generation, against 35\% for its
predecessor {[}42{]}. & Jul 2024 & demonstrated \\
Sparsity & Hardware-aligned, natively trainable sparse attention with
reported speedups at 64 thousand tokens of context {[}43{]}. & Feb 2025
& demonstrated \\
Decoding & Speculative decoding leaves the output distribution unchanged
and was measured at 2 to 3 times faster on a large translation model
{[}23{]}. & 30 Nov 2022 & established by proof in the source \\
Scaling & When lifetime inference demand is counted, the cost-optimal
model is smaller and trained on more tokens than the training-only rule
gives {[}44, 45{]}. & Jan 2024 & demonstrated \\
Training cost & A frontier-scale model reported at 2.788 million
accelerator-hours for full training, using latent attention and 8-bit
training {[}46{]}. & 26 Dec 2024 & demonstrated \\
Conformal inference & Finite-sample, distribution-free coverage for any
predictor, with proofs collected in a book {[}47{]}; a survey first
posted in October 2024 and published in 2025 covers structured,
unstructured and dynamic data {[}48{]}. & 18 Nov 2024 & established by
proof in the source \\
Inference with predictions & Valid intervals from a small labeled set
and many machine predictions, for any predictor {[}49{]}. & 2023 &
established by proof in the source \\
Sequential inference & E-processes and confidence sequences that stay
valid when a test is stopped at any time {[}50{]}. & Nov 2023 &
established by proof in the source \\
Causal estimation & Debiased machine learning gives valid inference on a
target parameter with flexible nuisance models, extended to identifiable
effects beyond back-door adjustment {[}51, 52{]}. & 2018; 2021 &
established by proof in the source \\
Compilers & Compiler passes for higher-order automatic differentiation
report up to 1.3 times on second-order operators, preliminary {[}53{]};
a 2026 methodology paper states that compiled JAX far exceeds Stan's
efficiency on modern hardware, with no benchmark numbers in its abstract
{[}54{]}. & 2025 & demonstrated \\
Probabilistic programming & A compiler that keeps probabilistic
semantics reports a geometric-mean 6.4 times over Stan and 53 times over
Turing on standard benchmarks {[}55{]}. & 2025 to 2026 & demonstrated \\
Energy measurement & An open benchmark measures inference energy across
40 architectures and reports time and energy Pareto configurations with
savings above 40\% in some cases {[}56{]}; one provider reports 0.24
watt-hours for a median text prompt in production {[}57{]}. & 2025 &
demonstrated \\
\end{longtable}
}
}

Two rates give the background. Compute needed to reach a fixed
language-model score has halved about every 8 months, with a 95\%
interval of about 5 to 14 months {[}58{]}. The price of reaching a fixed
benchmark level fell between 9 and 900 times a year across six
benchmarks, median about 50, with the fastest rates resting on under a
year of data {[}59{]}. Both are rates of change in the capability
frontier \(C_f(t)\) or a price, and neither measures the gap to
deployment.

Two limits apply to the table. Several rows are speedups against a
baseline chosen by the authors, so the numbers do not add or compare
across rows. And ``possible'' in these rows means demonstrated by the
authors in a stated setting, not available in a product.

\section{3. What the gap looks like in
data}\label{what-the-gap-looks-like-in-data}

\textbf{Long-run diffusion.} Across 15 technologies and 166 countries
from 1820 to 2003, the average lag between invention and adoption in a
country was about 45 years, and lags shrank over the period {[}60{]}.
The same history carries a warning about what adoption means. Electric
motors raised factory productivity only after plants were redesigned
around one motor per machine, a change that took decades {[}61{]}. A
later model attributes the dip and rise in measured productivity to
unmeasured complementary investment in processes and skills {[}62{]}. In
both accounts the delay is mostly the cost of reorganizing around the
tool, which in the notation above is \(K\).

\textbf{Open release to consumer hardware.} For models that fit one
consumer graphics card, the average lag behind the best model overall
was 6 to 12 months across four measures in 2025, with the largest lag,
12.4 months, on a preference score {[}63{]}. A broader analysis of
open-weight models of any size found about one year, with a range of 5
to 22 months by benchmark, and a later index-based analysis put the
average near 3 months {[}64{]}. These figures measure the distance
between two public frontiers. They are lower bounds on a lag to
products, because an open model on a card is not yet a product anyone is
paid to support.

\textbf{From paper to product.} Speculative decoding was posted in
November 2022. Its authors' employer reported applying it in products by
December 2024 {[}24{]}. A practitioner account from the same month says
that production use needed extra work on batching, quality and
reliability and helps mostly at limited batch sizes {[}65{]}. Two more
pairs are in Section 13. A few dated pairs are not a distribution, and
the protocol in Section 16 specifies how to collect many. A longer-run
measure of the paper-to-patent distance finds that patent citations to
articles are on average 14.4 years older than the patent on the front
page and 16.6 years when in-text citations are counted {[}66{]}, and
linked patents and papers are typically two to four citation steps apart
{[}67{]}. Those figures are for science broadly and are not about
machine learning in particular.

\textbf{Practice of 4-bit models.} Practitioner accounts date the point
at which 4-bit local inference became workable on consumer machines to
2023 {[}68{]}. The accounts are narrative, so this is opinion-class
evidence for a date and not a measurement.

\textbf{Firm adoption.} The United States Census Bureau's 2026
supplement found 18\% of firms using AI in a business function in
November 2025 to January 2026, 22\% expecting to within six months, and
57\% of users in three or fewer functions {[}34{]}. In the European
Union, 20.0\% of enterprises with 10 or more employees used at least one
AI technology in 2025, up 6.5 points from 2024, with 55.0\% of large,
30.4\% of medium and 17.0\% of small enterprises {[}33{]}. The OECD
reports 13.9\% adoption in 2024 and 39\% of large firms against 12\% of
small firms {[}35{]}. Chinese hiring data show generative-AI hiring at
21\% of listed firms against 0.2\% of unlisted ones at the end of 2025
{[}69{]}. Survey wording and sampling frames differ by source, and the
Census question was revised in November 2025, so rates are not
comparable across rows.

\textbf{Official statistics in the original languages.} Germany's
Statistisches Bundesamt reports that 20\% of firms with 10 or more
employees used AI in 2024, against 12\% in 2023 and 11\% in 2021, with
48\% of large, 28\% of medium and 17\% of small firms, and names lack of
knowledge and legal uncertainty as the most common reasons for not using
it {[}70{]}. Its 2025 table gives 26\% overall and 57\% for firms with
more than 250 employees, on a changed enterprise definition {[}31{]}.
Insee reports 10\% of firms in France using AI in 2024, four points
above 2023, 33\% of firms with 250 or more employees, and 42\% in
information and communication against 5\% or less in transport,
accommodation and food, and construction {[}37{]}. Spain's Instituto
Nacional de Estadística reports 21.1\% of firms with 10 or more
employees using AI in the first quarter of 2025, 8.7 points higher than
a year earlier {[}32{]}. Brazil's Cetic.br reports 13\% of firms using
AI in 2024, the same share as in 2023, concentrated in large firms and
in information and communication {[}36{]}. Korean ministry statistics,
as reported in the Korean press, give 32.9\% of firms using AI
technology or services in 2024 against 30.3\% in 2023, and 66.9\% for
firms with 250 or more workers {[}30{]}. Japan's white paper surveyed
workers in four countries and found any business use of generative AI at
55.2\% in Japan, 90.6\% in the United States, 90.3\% in Germany and
95.8\% in China. Among Japanese firms, 49.7\% had an active or
limited-use policy in 2024 against 42.7\% a year earlier, and the most
common Japanese concern, at 30.1\%, was not knowing how to use it
effectively, against 9.7\% in the United States and 10.4\% in China
{[}71{]}. China's CNNIC counts 515 million generative-AI users in June
2025, 266 million more than in December 2024, a penetration of 36.5\%
{[}72{]}. These are measurements of use, not of release. Within each
country the rate rises with firm size, which is consistent with a fixed
adoption cost \(K\) weighing less on a large budget, though the surveys
cannot separate that from other causes.

\textbf{Use at the level of people and tasks.} Nearly 40\% of United
States adults aged 18 to 64 used generative AI in late 2024, and about
23\% of workers had used it at work in the prior week {[}73{]}. A field
study of 5,179 customer-support agents found 14\% more resolutions per
hour on average and 34\% more for novice and low-skilled workers
{[}74{]}. A randomized trial with 16 experienced developers and 246
tasks found tasks with AI allowed took 19\% longer in that setting
{[}75{]}. These results differ because tasks, people and tools differ.
They show that deployment is not a single event, which is why \(D(t)\)
is defined from offers and not from use.

\textbf{What slows research itself.} Research effort has risen while
research productivity has fallen in most sectors, and keeping chip
density doubling needed more than 18 times the researchers of the early
1970s {[}76{]}. This bears on \(C_f(t)\), not on \(L\), and is noted so
the article does not read the capability frontier as free.

\section{4. Calendars: reporting and funding
periods}\label{calendars-reporting-and-funding-periods}

Decision dates in listed firms follow the reporting calendar. In a
survey of 401 financial executives, 78\% would give up economic value to
smooth reported earnings, and 55\% would avoid starting a very positive
net-present-value project if it meant missing the quarter's consensus
{[}77{]}. Comparable public firms invest less, and respond less to
investment opportunities, than private firms {[}78{]}. All of these are
demonstrated for earlier technologies and periods, and none is a
measurement of AI releases.

Attention follows the same calendar. The count of United States
large-company earnings calls citing ``AI'' reached 306 in the third
quarter of 2025, against a ten-year average of 86 {[}79{]}.
Earnings-call language has been used as a measurement instrument for
other topics, validated against manual reading {[}80{]}, and the same
instrument can be pointed at release topics.

Funding follows a longer calendar. A venture fund has a finite life,
typically ten years, and firms raise new funds every three to five years
{[}81{]}. Partnership agreements specify finite lifetimes, usually ten
years {[}82{]}. The two calendars differ in length, about a quarter and
three to five years, and both quantize when a commitment can be made.
Proposition 1 below uses only that structure.


\begin{center}
\openiefig{fig/latency-fig01-cadence-vs-quarter.png}{0.94}{https://research.openie.dev/figures/latency/cadence-vs-quarter.png}
\end{center}


\emph{Figure 1. Cadence against readiness: how a calendar adds wait
(model illustration). Quarterly decision dates with \(T = 0.25\) year;
readiness at a uniform random phase waits \(T/2\) on average. The
readiness phases are modeled, not observed.}

\section{5. Control and incumbent
response}\label{control-and-incumbent-response}

Two older results frame the incentive to deploy. Arrow's analysis finds
the incentive to invent smaller under monopoly, meaning barriers to
entry, than under competition, and counts a temporary monopoly that does
not bar entry as nearly competitive {[}83{]}. A later review of the
competition and innovation literature notes that preemption incentives
are a force missing from simple static models {[}84{]}. Incumbents also
miss changes that rearrange known components, because their routines
filter what they notice {[}85{]}. Neither is a statement about any firm
now.

Control over distribution is documented in public court records. In
September 2025 a United States district court barred exclusive
distribution contracts for several search and assistant products,
allowed non-exclusive default payments with terms generally limited to
one year, and ordered data sharing and syndication remedies {[}86{]}.
The article uses the record only as evidence that defaults and
distribution terms are priced and regulated, which is what the parameter
\(s\) in Section 6 represents.

\section{6. The model}\label{the-model}

\textbf{Setting.} One buyer population with demand \(Q(p) = A - p\) per
period for a task at price \(p\). An incumbent serves it at cost \(c_0\)
with an existing method. A public method becomes available to anyone at
cost \(c_e(t)\), non-increasing in time. A buyer can switch to a seller
using the public method at friction \(s\) per task, so the incumbent's
price ceiling is \(u(t) = c_e(t) + s\). The incumbent may release an
offering built on the public method at one-time cost \(K\). Decisions
are made only at dates \(kT\).

\textbf{Assumptions.}

\begin{itemize}
\tightlist
\item
  A1. Linear demand \(A - p\), with \(c_0 < A - s\).
\item
  A2. The public cost declines: \(c_e\) is non-increasing and continuous
  with \(c_e(0) = c_0\), and \(c_e(t) \ge h \ge 0\).
\item
  A3. The ceiling \(u\) binds on the incumbent's price, that is
  \(u < \tfrac12 (A + c)\) for the cost \(c\) it faces. This holds for
  the numbers in Section 7.
\item
  A4. The incumbent maximizes discounted profit at rate \(r\) subject to
  a reporting rule: a share \(\varphi \in [0,1]\) of \(K\) is expensed
  in the release period, and the release period's reported earnings must
  not fall short of the pre-release level by more than \(w\).
\item
  A5. After release the incumbent's cost is \(c_e(t)\), the cost of the
  public method.
\item
  A6. Flow gains from release are treated as constant after the release
  date (steady-state approximation).
\end{itemize}

\textbf{Proposition 1 (calendar quantization).} If a method becomes
ready at time \(t_r\) and the firm releases at the first decision date
at or after \(t_r\) when release is worthwhile, the wait is
\(w^{cad} = \lceil t_r / T \rceil T - t_r\). If \(t_r \bmod T\) is
uniform on \([0, T)\), then \(E[w^{cad}] = T/2\) and
\(\mathrm{Var}[w^{cad}] = T^2/12\).

\emph{Derivation.} Write \(t_r = mT + x\) with \(x\) uniform on
\([0,T)\). For \(x > 0\) the next date is \((m+1)T\) and the wait is
\(T - x\), which is uniform on \((0, T]\). A uniform variable on an
interval of length \(T\) has mean \(T/2\) and variance \(T^2/12\). The
case \(x = 0\) has probability zero. \(\square\)

\emph{Consequence.} A firm deciding on a quarterly calendar adds on
average about 6 weeks, and one deciding on a funding cycle of 3 to 5
years adds 1.5 to 2.5 years, to the readiness delay \(d\), with no
engineering cause. The assumption that readiness is uniform within a
period is the weak part, and Prediction 1 tests the converse.

\textbf{Proposition 2 (release threshold and forced release).} Define
the flow gain \(g(c) = (c_0 - c)(A - c - s)\) for the outside cost
\(c = c_e\). On \(0 \le c < \min(c_0, A - s)\), \(g\) is strictly
decreasing in \(c\). Without calendar or reporting constraints, release
is optimal exactly when \(g(c_e(t)) \ge rK\), that is when
\(c_e(t) \le c^\dagger\), where \(c^\dagger\) is the unique solution of
\(g(c^\dagger) = rK\) if it exists. If \(c_e \le c_0 - s\), holding
earns no profit, and release is preferred to holding. A reporting rule
can delay it but cannot make holding profitable.

\emph{Derivation.} Under A3 the incumbent holding at cost \(c_0\) prices
at the ceiling, so its profit is \((u - c_0)(A - u)\). After release its
cost is \(c\) and its profit is \((u - c)(A - u)\). The difference is
\((c_0 - c)(A - u) = (c_0 - c)(A - c - s) = g(c)\). Differentiating,
\(g'(c) = -(A - c - s) - (c_0 - c) < 0\) on the stated range, so \(g\)
is strictly decreasing. With A6 the value of release is \(g/r - K\),
which is non-negative exactly when \(g \ge rK\). Holding yields profit
\((u - c_0)(A - u)\), which is zero at \(u = c_0\), that is at
\(c_e = c_0 - s\), and negative below. Release at that point yields
profit \(s(A - c - s) > 0\) by A1. \(\square\)

\textbf{Proposition 3 (reporting gate, delay and a never-release
region).} Under A4 the release-period reported earnings change is
\(g \cdot T - \varphi K\). Release is allowed if this is at least
\(-w\), that is if \(g \ge (\varphi K - w)/T\). Define
\(G_{eff} = \max\{rK, (\varphi K - w)/T\}\) and
\(c_{eff} = g^{-1}(G_{eff})\). Then \(c_{eff} \le c^\dagger\), the
release date moves from \(t(c^\dagger)\) to \(t(c_{eff})\) or to the
forced-release point \(c_0 - s\), whichever is earlier, and for
\(c_e(t) = c_e(0) e^{-\lambda t}\) the delay is
\((1/\lambda)\ln(c^\dagger / c_{eff})\) when \(c_{eff} \ge c_0 - s\). If
\(G_{eff} > g(h)\), where \(h\) is the cost floor, release never occurs
on the NPV and gate conditions, and only the forced release at
\(c_0 - s\) applies if \(c_0 - s \ge h\).

\emph{Derivation.} Reported earnings in the release period gain \(gT\)
of operating profit and lose the expensed part \(\varphi K\) of the
adoption cost. Requiring the net to be at least \(-w\) gives the stated
inequality. Since \(g\) is strictly decreasing (Proposition 2), the
larger of two required gains corresponds to the smaller cost threshold,
so \(c_{eff} \le c^\dagger\). With \(c_e(t) = c_e(0) e^{-\lambda t}\)
the time to reach \(c\) is \(t(c) = (1/\lambda)\ln(c_e(0)/c)\), and the
difference of two such times is the stated delay. The attainable gain is
at most \(g(h)\) because \(c_e \ge h\) by A2. \(\square\)

\textbf{Proposition 4 (forcing function).} Suppose an automation tool
lowers the adoption cost from \(K\) to \(K'\). (a) The NPV threshold
\(c^\dagger\) rises, so release is optimal earlier. (b) If
\(\varphi K' \le w\), the reporting gate does not bind,
\(G_{eff} = rK'\), and \(w^{gate} = 0\). (c) The total latency of the
firm falls to \(d' + w^{cad}\), where \(d'\) is the technical readiness
delay with the tool.

\emph{Derivation.} (a) \(g(c^\dagger) = rK\) and \(g\) is strictly
decreasing, so lowering \(K\) lowers the required gain and raises
\(c^\dagger\). (b) The gate condition \(g \ge (\varphi K' - w)/T\) is
satisfied by every \(g \ge 0\) when the right side is non-positive,
which is \(\varphi K' \le w\). Then \(G_{eff} = rK'\). (c) Follows from
Definition 2 and Proposition 3. \(\square\)

\emph{Reading.} A tool acts as a forcing function in two separate ways.
It lowers \(K\) and so removes the gate delay once \(K\) passes under
\(w/\varphi\). And it lowers \(c_e\) for every seller, which forces
release at \(c_e \le c_0 - s\), where holding earns nothing. The first
is a property of the adopting firm and the second a property of the
market.

\textbf{Proposition 5 (market latency is set by the least constrained
firm).} If \(D(t) = \min_i c_i(t)\) over firms that can sell to the
buyer population, then \(L(c) = \min_i L_i(c)\).

\emph{Derivation.}
\(t_D(c) = \inf\{t : \min_i c_i(t) \le c\} = \min_i t_i(c)\), and
\(t_C(c)\) does not depend on \(i\). Subtracting gives the result.
\(\square\)

\emph{Consequence.} A calendar or gate that delays one firm does not
delay the market when another firm can reach the same buyers without
that delay. The latency observed in a market is then the latency of the
fastest firm that has access. Where access is restricted, because buyers
cannot reach other sellers, the minimum is taken over a smaller set. In
the model this is the role of \(s\).

\textbf{Proposition 6 (what remains to charge for).} After release, with
the ceiling binding, price is \(p = c_e + s\), profit is
\(\pi = s(A - c_e - s)\), and the margin per task is \(s\), independent
of \(c_e\). Price changes one-for-one with cost, while the price to cost
ratio \(p/c_e = 1 + s/c_e\) rises as cost falls. The marginal value of
friction is \(\partial \pi / \partial s = A - c_e - 2s\), which is
positive for \(s < (A - c_e)/2\) and increases as \(c_e\) falls.

\emph{Derivation.} Under A3 the incumbent charges the ceiling
\(u = c_e + s\). With cost \(c_e\), the margin is \(u - c_e = s\) and
volume is \(A - u\), giving \(\pi = s(A - c_e - s)\). The ratio follows
from \(p/c_e = (c_e + s)/c_e\). Differentiating \(\pi\) in \(s\) gives
\(A - c_e - 2s\), and
\(\partial^2\pi / \partial s\,\partial(-c_e) = 1 > 0\). \(\square\)

\emph{Reading.} As the cost of computing falls, the part of price that
reflects computing shrinks, and the remaining margin depends on how
costly it is for a buyer to leave. Spending that raises \(s\) becomes
worth more. In the model this is a statement about incentives and not a
judgment of any firm. Policies that lower \(s\), such as limits on
default terms or portable data formats, lower the margin by the same
amount, and the model does not say whether that is desirable.

\section{7. A numerical illustration, not a
fit}\label{a-numerical-illustration-not-a-fit}

The script \texttt{model\_check.py} in the work folder reproduces these
numbers with the standard library. Parameters: \(A = 10\), \(c_0 = 4\),
\(s = 2\), \(r = 0.10\) per year, \(T = 0.25\) year, \(\varphi = 1\),
\(w = 0.1\), and outside cost halving every 8 months so
\(c_e(t) = 4 e^{-\lambda t}\). The 8-month rate echoes {[}58{]} and is a
choice for illustration, not an estimate for cost.

\textbf{Table 2. Release date with and without an adoption-cost-lowering
tool.}

{\def\LTcaptype{none} % do not increment counter
{\scriptsize\setlength{\tabcolsep}{3.0pt}\renewcommand{\arraystretch}{1.14}
\begin{longtable}[]{@{}
  >{\raggedright\arraybackslash}p{(\linewidth - 10\tabcolsep) * \real{0.2500}}
  >{\raggedright\arraybackslash}p{(\linewidth - 10\tabcolsep) * \real{0.0729}}
  >{\raggedright\arraybackslash}p{(\linewidth - 10\tabcolsep) * \real{0.1979}}
  >{\raggedright\arraybackslash}p{(\linewidth - 10\tabcolsep) * \real{0.1771}}
  >{\raggedright\arraybackslash}p{(\linewidth - 10\tabcolsep) * \real{0.1458}}
  >{\raggedright\arraybackslash}p{(\linewidth - 10\tabcolsep) * \real{0.1562}}@{}}
\toprule\noalign{}
\begin{minipage}[b]{\linewidth}\raggedright
\textbf{Case}
\end{minipage} & \begin{minipage}[b]{\linewidth}\raggedright
\textbf{\(K\)}
\end{minipage} & \begin{minipage}[b]{\linewidth}\raggedright
\textbf{NPV threshold}
\end{minipage} & \begin{minipage}[b]{\linewidth}\raggedright
\textbf{Gate threshold}
\end{minipage} & \begin{minipage}[b]{\linewidth}\raggedright
\textbf{Release at}
\end{minipage} & \begin{minipage}[b]{\linewidth}\raggedright
\textbf{Gate delay}
\end{minipage} \\
\midrule\noalign{}
\endhead
\bottomrule\noalign{}
\endlastfoot
Before the tool & 2.0 & \(c = 3.951\), day 4 & \(c = 2.594\) & month
5.00 & 4.85 months \\
Tool lowers adoption cost & 0.1 & \(c = 3.998\) & \(c = 3.998\) & day
0.2 & 0 \\
\end{longtable}
}
}

In the first row the gate requires a gain of
\((\varphi K - w)/T = 7.6\), far above the NPV requirement \(rK = 0.2\),
so the incentive-delay part \(w^{gate}\) is 4.85 months. That is longer
than the calendar term of Proposition 1 for a quarterly firm, which is
1.5 months on average. In the second row
\(\varphi K' = 0.1 \le w = 0.1\), so the gate does not bind and only the
NPV condition remains, as in Proposition 4(b).


\begin{center}
\openiefig{fig/latency-fig02-forcing-function-threshold.png}{0.94}{https://research.openie.dev/figures/latency/forcing-function-threshold.png}
\end{center}


\emph{Figure 2. The forcing function: a lower adoption cost removes the
gate (model illustration). Gain \(g(c)\) against the two thresholds,
with the parameters of this section. Data: \texttt{model\_check.py} and
Table 2.}

\textbf{Table 3. After release, with \(s = 2\).}

{\def\LTcaptype{none} % do not increment counter
{\footnotesize\setlength{\tabcolsep}{3.8pt}\renewcommand{\arraystretch}{1.15}
\begin{longtable}[]{@{}
  >{\raggedright\arraybackslash}p{(\linewidth - 8\tabcolsep) * \real{0.1146}}
  >{\raggedright\arraybackslash}p{(\linewidth - 8\tabcolsep) * \real{0.1875}}
  >{\raggedright\arraybackslash}p{(\linewidth - 8\tabcolsep) * \real{0.1875}}
  >{\raggedright\arraybackslash}p{(\linewidth - 8\tabcolsep) * \real{0.1458}}
  >{\raggedright\arraybackslash}p{(\linewidth - 8\tabcolsep) * \real{0.3646}}@{}}
\toprule\noalign{}
\begin{minipage}[b]{\linewidth}\raggedright
\textbf{\(c_e\)}
\end{minipage} & \begin{minipage}[b]{\linewidth}\raggedright
\textbf{Price \(p\)}
\end{minipage} & \begin{minipage}[b]{\linewidth}\raggedright
\textbf{\(p / c_e\)}
\end{minipage} & \begin{minipage}[b]{\linewidth}\raggedright
\textbf{Profit}
\end{minipage} & \begin{minipage}[b]{\linewidth}\raggedright
\textbf{\(\partial \pi / \partial s\)}
\end{minipage} \\
\midrule\noalign{}
\endhead
\bottomrule\noalign{}
\endlastfoot
4.0 & 6.00 & 1.50 & 8.00 & 2.00 \\
3.0 & 5.00 & 1.67 & 10.00 & 3.00 \\
2.0 & 4.00 & 2.00 & 12.00 & 4.00 \\
1.0 & 3.00 & 3.00 & 14.00 & 5.00 \\
0.5 & 2.50 & 5.00 & 15.00 & 5.50 \\
\end{longtable}
}
}

Profit rises as cost falls because linear demand lets volume expand.
With satiated demand the profit column would flatten, and the claim
about margin staying at \(s\) would not change. A simulation of
Proposition 1 with \(T = 0.25\) and uniform readiness gives a mean wait
of 0.1250 against \(T/2 = 0.125\) and variance 0.00521 against
\(T^2/12 = 0.00521\) (200,000 draws, seed 7).

\section{8. Predictions and what would refute
them}\label{predictions-and-what-would-refute-them}

Each prediction names the data, the test and the result that would count
against the model. Predictions are stated before any of the data was
collected.

\textbf{Prediction 1 (calendar clustering).} The dates on which firms
with quarterly reporting make previously open methods generally
available are not uniform within the quarter. \emph{Test:} for a sample
of at least 200 releases from listed firms, compute each date's phase
within the firm's fiscal quarter and apply Rayleigh and Kuiper tests
against uniformity. \emph{Refuted if} the null of uniform phase is not
rejected at 5\% with 200 or more independent items.

\textbf{Prediction 2 (listed against unlisted).} For the same method and
comparable difficulty, the median delay from open release to a generally
available product is longer for listed firms than for firms funded
outside a reporting calendar. \emph{Test:} matched pairs by method, with
firm type as the treatment and hazard ratios for time to release.
\emph{Refuted if} the 95\% interval of the hazard ratio (listed over
unlisted) lies entirely above 0.8.

\textbf{Prediction 3 (tools lower the gate).} After an automation tool
that lowers porting cost becomes available for a hardware or software
target, time from open release to general availability on that target
falls. \emph{Test:} difference in differences with targets served and
not served by the tool, with release delay in days as outcome.
\emph{Refuted if} the 95\% interval of the estimated change lies
entirely above \(-30\) days. The 30-day margin is a protocol choice made
before data collection.

\textbf{Prediction 4 (price outruns cost).} Where joule and hardware
costs per task are measured, the ratio of list price to cost per task
rises as measured cost falls. \emph{Test:} panel of services with a
stated task, using open energy measurements and list prices.
\emph{Refuted if} the ratio is flat or falls in more than half of tasks
over a period in which cost fell by a factor of 2 or more.

\textbf{Prediction 5 (spending moves to friction).} In a market where
cost per task falls, the share of seller revenue spent on distribution,
defaults and bundling rises. \emph{Test:} disclosed distribution
payments and bundling terms in filings and court records against
revenue. \emph{Refuted if} the share is flat or falls while cost per
task falls by a factor of 2 or more.

\textbf{Prediction 6 (market latency follows the fastest firm).} In
markets where buyers can reach more than one seller at low friction,
market latency is close to the minimum firm latency, and in markets with
high measured friction it is closer to the median. \emph{Test:} regress
market latency on measured switching friction. \emph{Refuted if} the
95\% interval of the slope on friction lies entirely at or below zero.

\textbf{Predictions on the financing and cost of the build-out.} The
predictions in Sections 10 and 12 concern the money flows and the cost
of an accelerator-hour. They are stated in those sections, and they are
tested on filings and listings dated after this article.

\textbf{What would refute the model itself.} A finding that release
dates are uniform within the quarter for listed firms (Prediction 1),
that listed and unlisted firms release at the same pace (Prediction 2)
and that porting-tool arrival does not change delay (Prediction 3) would
leave no role for calendars or reporting gates, and the artificial
latency of Definition 2 would be zero in measurement. The model would
then reduce to engineering delay \(d\).

\section{9. The money loop: what the filings
document}\label{the-money-loop-what-the-filings-document}

A common description of the current build-out is a circular economy:
labs buy compute from clouds, clouds and neoclouds buy chips and memory,
and chip makers and clouds invest in the labs. This section sets out
what dated documents say about that description. It gives no verdict.
Each item is marked \textbf{documented} when a filing, a dated release
or an official statistic states the amount and the counterparties,
\textbf{inferred} when it follows by arithmetic from documented figures
under a stated assumption, and \textbf{not found} when the sources read
do not say.


\begin{center}
\openiefig{fig/latency-fig03-capital-compute-flow.png}{0.94}{https://research.openie.dev/figures/latency/capital-compute-flow.png}
\end{center}


\emph{Figure 3. Capital and compute commitments among labs, clouds, chip
makers and finance. Arrows follow the direction of money or of the
purchase commitment and are not netted. Data:
\texttt{scripts/\allowbreak{}figures/\allowbreak{}latency\_\allowbreak{}figures.\allowbreak{}py} with
\texttt{latency\_\allowbreak{}data/\allowbreak{}flows.\allowbreak{}csv}; sources are listed in Table 4.}

\textbf{Capital supplied to customers.} In the round of 122 billion
dollars that OpenAI closed in March 2026, NVIDIA is reported at 30
billion dollars, Amazon at up to 50 billion and SoftBank at 30 billion
{[}9{]}. An earlier letter of intent had said NVIDIA intended to invest
up to 100 billion dollars progressively as each gigawatt is deployed,
which is a statement of intent and not a closed transaction {[}87{]}.
Microsoft reports about 27\% of the OpenAI public benefit corporation
and 13 billion dollars of total funding commitments accounted under the
equity method {[}7{]}. For Anthropic, NVIDIA and Microsoft committed to
invest up to 10 and 5 billion dollars {[}10{]}. NVIDIA bought 2 billion
dollars of CoreWeave Class A stock after the end of that company's 2025
fiscal year {[}4{]}. AMD issued OpenAI a warrant for up to 160 million
shares at 0.01 dollars per share, vesting in tranches tied to purchases
of AMD Instinct GPU products {[}88{]}.

\textbf{Purchases committed by customers.} OpenAI contracted to buy an
incremental 250 billion dollars of Azure services {[}7, 89{]}, a 38
billion dollar agreement with AWS that the announcement describes as
growing over seven years {[}8{]}, and up to 6.5 billion dollars from
CoreWeave, while Meta has an order form with CoreWeave of up to 14.2
billion {[}4{]}. Anthropic committed to buy 30 billion dollars of Azure
compute and to contract up to one gigawatt more {[}10{]}. NVIDIA itself
reports 27 billion dollars of multi-year cloud service commitments and
3.5 billion dollars of land, power and shell guarantees to early-stage
partners {[}3{]}.

\textbf{Table 4. What is documented, what is inferred, and what the
sources do not say.}

{\def\LTcaptype{none} % do not increment counter
{\footnotesize\setlength{\tabcolsep}{3.8pt}\renewcommand{\arraystretch}{1.15}
\begin{longtable}[]{@{}
  >{\raggedright\arraybackslash}p{(\linewidth - 6\tabcolsep) * \real{0.1453}}
  >{\raggedright\arraybackslash}p{(\linewidth - 6\tabcolsep) * \real{0.3023}}
  >{\raggedright\arraybackslash}p{(\linewidth - 6\tabcolsep) * \real{0.2733}}
  >{\raggedright\arraybackslash}p{(\linewidth - 6\tabcolsep) * \real{0.2791}}@{}}
\toprule\noalign{}
\begin{minipage}[b]{\linewidth}\raggedright
\textbf{Item}
\end{minipage} & \begin{minipage}[b]{\linewidth}\raggedright
\textbf{Documented}
\end{minipage} & \begin{minipage}[b]{\linewidth}\raggedright
\textbf{Inferred, with the assumption}
\end{minipage} & \begin{minipage}[b]{\linewidth}\raggedright
\textbf{Not found in the sources read}
\end{minipage} \\
\midrule\noalign{}
\endhead
\bottomrule\noalign{}
\endlastfoot
Equity from suppliers into customers & The amounts and dates above {[}9,
10, 4{]} & None drawn & Whether any equity is conditioned on purchases,
except the AMD warrant, whose vesting is tied to GPU purchases in the
filing {[}88{]} \\
Purchases by labs from clouds & Azure 250 and AWS 38 billion dollars for
OpenAI; Azure 30 billion for Anthropic {[}7, 8, 10{]} & the AWS figure
averages about 5.4 billion dollars a year if spread evenly over seven
years & The term and ramp of the Azure commitment; the conditions
attached to the 50 billion dollar Amazon figure \\
Customer concentration & NVIDIA direct customers at 22\% and 14\% of
revenue, with system builders and distributors among direct customers
{[}3{]}; CoreWeave with Microsoft at 67\% of 2025 revenue {[}4{]} & A
lower bound on the Herfindahl index from the disclosed shares: at least
0.068 for NVIDIA and 0.449 for CoreWeave & Who stands behind NVIDIA's
direct customers \\
Customer prepayments & Oracle's deferred revenue from customer
prepayments with a significant financing component rose by 11.4 billion
dollars in the quarter, inside operating cash flow of 23.1 billion
{[}5{]} & About half of that quarter's operating cash flow came from
prepayments & Which customers prepaid \\
Venture-level financing & Blue Owl funds hold 80\% and Meta 20\% of a
venture for the Hyperion campus; the funds contributed 7.0 billion
dollars cash against about 27 billion dollars of development cost {[}2,
1{]} & The cash is about a quarter of development cost & Lease terms,
debt raised by the venture and guarantees \\
Depreciation lives & Stated lives for servers in Figure 12 {[}4, 14, 15,
16, 17{]} & Moving a life from 5 to 6 years lowers the straight-line
annual charge on the same assets by one sixth & The service life that
GPUs reach in use \\
Cash capital against operating cash & The annual series in Figure 8
{[}11{]} & The ratio of the two & Finance leases and off-balance-sheet
leases \\
\end{longtable}
}
}

The term \textbf{customer} in this table follows each filing's own use.
Money is fungible, so a filing can show that capital went to a customer
and that the customer paid the supplier, and it cannot show that the one
funded the other. Section 10 states a test that works within that limit.


\begin{center}
\openiefig{fig/latency-fig04-hyperion-structure.png}{0.94}{https://research.openie.dev/figures/latency/hyperion-structure.png}
\end{center}


\emph{Figure 4. One financing structure as filed: the Hyperion campus
venture. USD billions as printed. Data: Meta press release and quarterly
report {[}1, 2{]}.}

\textbf{Concentration and backlog.} Disclosed customer shares are high
at the leveraged nodes and lower at the largest chip maker. CoreWeave
reports its top three customers at 73\% of 2023 revenue, its top two at
77\% in 2024 and Microsoft alone at 67\% in 2025 {[}4{]}. Its remaining
performance obligations of 60.7 billion dollars are 11.8 times its 2025
revenue of 5.131 billion. Oracle reports obligations of 664 billion
dollars at 31 August 2026 and expects about 13\% of them to be
recognized in the next twelve months {[}5, 6{]}, against fiscal 2026
revenue of 67.4 billion {[}11{]}. OpenAI reported revenue of 13.1
billion dollars for 2025 and stated a rate of 2 billion dollars a month
in March 2026 {[}9{]}. Contracted purchases are multi-year and
conditional on terms that these figures do not show, so they are not
revenue in one year and are not additive with the figures above.


\begin{center}
\openiefig{fig/latency-fig05-customer-concentration.png}{0.94}{https://research.openie.dev/figures/latency/customer-concentration.png}
\end{center}


\emph{Figure 5. Customer concentration as disclosed. Share of revenue
stated in annual reports; customers are named only where the filing
names them. Data: annual reports {[}3, 4{]}.}


\begin{center}
\openiefig{fig/latency-fig06-commitments-vs-revenue.png}{0.94}{https://research.openie.dev/figures/latency/commitments-vs-revenue.png}
\end{center}


\emph{Figure 6. Contracted commitments against revenue earned, USD
billions on a log scale. Commitments are multi-year and conditional;
revenue is one fiscal year. Data: filings and releases {[}4, 5, 6, 7, 8,
9{]}.}

\textbf{Labs are also buying elsewhere.} Anthropic announced a plan to
use up to one million TPUs, worth tens of billions of dollars and well
over a gigawatt of capacity in 2026 {[}90{]}. OpenAI announced a plan to
deploy 10 gigawatts of accelerators of its own design {[}91{]} and holds
the AMD warrant above. A lab that buys from several suppliers and
designs its own parts is not a single customer of any one of them, which
matters for the bound in the next section.

\section{10. A closed-loop test and its
predictions}\label{a-closed-loop-test-and-its-predictions}

The question the documents cannot answer by themselves is how much of a
supplier's revenue is paid for with money the supplier supplied. The
test below turns it into quantities that a filing can bound.

For a window \(W\) (a fiscal year), let \(F_{sc}\) be capital supplied
by supplier \(s\) to customer \(c\) (equity, loans, guarantees, warrants
valued at issue), let \(R_{cs}\) be the revenue \(s\) receives from
\(c\), let \(R_s\) be the total revenue of \(s\), and let \(Z_c\) be the
capital \(c\) raises from all sources.

\textbf{Definition 3 (round-trip ratio).}
\(\rho_{sc} = F_{sc} / R_{cs}\), capital supplied per dollar of revenue
received from the same counterparty.

\textbf{Definition 4 (closed-loop share bound).}
\(\bar\varphi_s = \sum_c \min(F_{sc}, R_{cs}) / R_s\). A customer cannot
have paid \(s\) with more of \(s\)'s own capital than \(s\) supplied,
nor than it paid, so \(\bar\varphi_s\) is an upper bound on the share of
\(R_s\) that could have been funded by \(s\)'s own capital. The
customer-side counterpart is \(\psi_c = \sum_s F_{sc} / Z_c\), the share
of capital raised by \(c\) that came from its own suppliers. Both are
bounds computed from disclosed figures. They are not estimates of what
happened, because money is fungible.

\textbf{Proposition 7 (when a supplier-financed purchase pays).} Let
supplier \(s\) sell at price \(p\) with unit cost \(u\) and gross margin
\(m = (p - u)/p\). Let it supply capital \(F\) that induces a share
\(\theta\) of sales \(Q\) that would not otherwise occur, and let it
recover value \(V\) from its stake. The incremental profit is
\(\pi = \theta (p - u) Q + V - F\). If all of \(F\) is spent with \(s\),
so that \(pQ = F\), then \(\pi \ge 0\) if and only if
\(\theta \ge \theta^* = (1 - v)/m\), where \(v = V/F\).

\emph{Derivation.} Substitute \((p-u)Q = m\,pQ = mF\) and solve
\(\theta m F + vF - F \ge 0\). \(\blacksquare\) (By hand, not
machine-checked.)

\emph{Consequence.} Since \(\theta \le 1\), supplier financing pays
through induced sales alone only when \(v \ge 1 - m\). With the gross
margin of 71.1\% reported by NVIDIA for fiscal 2026 {[}3{]}, the stake
must be expected to recover at least 28.9\% of the capital. Model output
for the breakeven induced share \(\theta^*\) at recovery 0, 25\%, 50\%,
75\% and 100\% is 1.41, 1.05, 0.70, 0.35 and 0, so with no recovery no
share of induced sales is enough.


\begin{center}
\openiefig{fig/latency-fig07-closed-loop-bounds.png}{0.94}{https://research.openie.dev/figures/latency/closed-loop-bounds.png}
\end{center}


\emph{Figure 7. Closed-loop test, upper bounds from documented figures
(model output). NVIDIA alone is 30 of the 122 billion dollar OpenAI
round; adding Amazon at its stated maximum gives 80 of 122; directed
capital announced from November 2025 to March 2026 (30, 10 and 2 billion
dollars) is compared with NVIDIA's fiscal 2026 revenue of 215.9 billion
dollars. Data: \texttt{latency\_\allowbreak{}data/\allowbreak{}model\_\allowbreak{}outputs.\allowbreak{}json}, computed by
\texttt{model\_\allowbreak{}airace.\allowbreak{}py} from {[}9, 10, 4, 3{]}.}

\textbf{Bounds from the documents.} NVIDIA's share of the OpenAI round
is at most 24.6\%, and with Amazon at its stated maximum the two reach
at most 65.6\%, with Microsoft's part not disclosed (model output).
Capital that NVIDIA announced or was reported to supply to customers
from November 2025 to March 2026 is 42 billion dollars (30 for OpenAI,
up to 10 for Anthropic and 2 for CoreWeave), which is 19.5\% of its
fiscal 2026 revenue of 215.9 billion and 41\% of its operating cash flow
of 102.7 billion {[}3, 11{]}. The first figure is a scale comparison,
since the amounts and the revenue year differ. Read as a bound, it means
that if every one of those dollars came back as NVIDIA revenue within
one year, at most about a fifth of that year's revenue would be funded
by NVIDIA's own capital.

\textbf{Predictions.} Each names data, a test and a result that would
count against the claim. These predictions were written after the
filings above were read, and they are tested on filings dated after this
article.

\textbf{Prediction 7 (the loop share stays bounded).} For each chip or
cloud supplier with disclosed directed capital, \(\bar\varphi_s\)
computed from each annual report stays below 0.25 while revenue from
customers that received no directed capital grows at least as fast as
revenue from those that did. \emph{Test:} annual filings, with customers
named where filings name them. \emph{Refuted if} \(\bar\varphi_s\)
exceeds 0.5 in a fiscal year, or if revenue growth is concentrated in
customers that received directed capital. Values between 0.25 and 0.5
are inconclusive.

\textbf{Prediction 8 (end-user revenue outgrows commitments).} For a
lab, the ratio of contracted purchases to trailing twelve-month revenue
falls over four consecutive quarters while revenue rises. \emph{Test:}
lab statements and cloud backlog disclosures. \emph{Refuted if} the
ratio rises for four consecutive quarters.

\textbf{Prediction 9 (stress shows first where concentration and
leverage are highest).} When a lab revises a compute commitment
downward, the size of revisions to guidance at supplier nodes is ordered
by their disclosed top-customer share. \emph{Test:} guidance revisions
after any such event, ordered by disclosed top-customer share.
\emph{Refuted if} the order is unrelated to concentration.

\textbf{Prediction 10 (substitution reduces committed capacity per
dollar of revenue).} Where open methods and cheaper hardware reduce
accelerator time per task (Section 12), committed gigawatts per dollar
of lab revenue fall. \emph{Test:} announced capacity and lab revenue by
quarter. \emph{Refuted if} committed gigawatts per dollar of revenue
rise for four consecutive quarters while measured accelerator time per
task also falls on open benchmarks.

\section{11. The strongest
counter-evidence}\label{the-strongest-counter-evidence}

The reading that the build-out is mostly a closed loop meets four sets
of documents that point the other way, and one set that points toward
fragility.

\textbf{Most spending is funded from operations at the largest buyers.}
In their latest fiscal years, cash capital spending was 0.63 of
operating cash flow at Microsoft, 0.94 at Amazon, 0.56 at Alphabet and
0.60 at Meta {[}11{]}. The spending is large, and it is paid for out of
cash that the same firms earn from other businesses.


\begin{center}
\openiefig{fig/latency-fig08-capex-vs-cash-flow.png}{0.94}{https://research.openie.dev/figures/latency/capex-vs-cash-flow.png}
\end{center}


\emph{Figure 8. Cash capital spending against cash from operations,
annual, USD billions. A ratio above 1 means spending exceeded operating
cash flow that year. Excludes finance leases and off-balance-sheet
leases. Data: SEC XBRL annual series in
\texttt{latency\_\allowbreak{}data/\allowbreak{}capex\_\allowbreak{}ocf.\allowbreak{}csv} {[}11{]}.}

\textbf{The largest chip supplier could fund its directed capital from
its own cash.} The 42 billion dollars above is 41\% of one year's
operating cash flow at NVIDIA {[}3, 11{]}, and its two largest direct
customers are 22\% and 14\% of revenue, with system builders and
distributors among them {[}3{]}.

\textbf{Revenue from end users is real and growing in the statements
read.} OpenAI reported 13.1 billion dollars of revenue for 2025 and a
rate of 2 billion dollars a month in March 2026, which annualizes to 24
billion (arithmetic) {[}9{]}.

\textbf{Labs diversify.} Anthropic's plan for up to one million TPUs
alongside its Azure and NVIDIA commitments, and OpenAI's plans for its
own accelerators and the AMD warrant, mean that the labs named here have
more than one chip source in their announced plans {[}90, 10, 91, 88{]}.

\textbf{What the counter-evidence does not remove.} Cash capital
spending was 1.74 times operating cash flow at Oracle in fiscal 2026 and
3.37 times at CoreWeave in 2025 {[}11{]}. Oracle's quarter ended 31
August 2026 shows operating cash flow of 23.1 billion dollars with 11.4
billion of that from customer prepayments, and free cash flow of
negative 5 billion {[}5, 6{]}. CoreWeave's obligations are 11.8 times
its revenue, and its customer base is concentrated {[}4{]}. These are
the nodes where a shortfall in a lab's funding would arrive first, which
is what Prediction 9 tests.

\textbf{Net reading, as far as the documents go.} Loops of tens of
billions of dollars are documented between named suppliers and
customers. The documents bound the share of the largest chip supplier's
revenue that its own capital could have funded under a fifth, and show
most hyperscaler spending paid from operating cash. The concentrated and
leveraged nodes, a cloud with prepayments and a neocloud with a single
dominant customer, are where the loop assumption would matter most.
Terms of private agreements are not public, so the documents give bounds
and not a measurement.

\section{12. What a task costs: the infrastructure cost
model}\label{what-a-task-costs-the-infrastructure-cost-model}

Section 0 defines cost per task as \(c = \kappa e + h\). This section
fills in \(\kappa\) and \(h\) from the cost of an accelerator-hour and
ties the result to \(C_f(t)\), \(D(t)\) and \(L\).

\textbf{Accelerator-hour.} Let \(K_{acc}\) be capital (written with a
subscript to keep it apart from the adoption cost \(K\) of Section 6)
per accelerator, including the share of facility, power and cooling,
with straight-line life \(y\) years and utilization \(u\). Capital per
utilized hour is \(a = K_{acc} / (8760\, y\, u)\). Let \(P\) be facility
power per accelerator in kilowatts and \(\kappa_e\) the price of
electricity per kilowatt-hour. A task that occupies \(\tau\) utilized
hours costs
\[c = \tau\,(a + \kappa_e P), \qquad e = 3.6\times 10^{6}\, P\,\tau \ \text{joules}.\]

\textbf{Proposition 8 (the capital-loaded price of a joule).} At fixed
hardware, \(c = \kappa^{*} e\) with \(\kappa^{*} = \kappa_e + a/P\) per
kilowatt-hour equivalent, so joules per task stays the right unit of
account, and the price that converts joules into dollars is dominated by
\(a/P\) whenever \(a/P \gg \kappa_e\). \emph{Derivation.} Eliminate
\(\tau\) between the two expressions. \(\blacksquare\) (By hand.)

\textbf{Inputs.} The Epoch AI data-center dataset gives, for 74 sites
with both figures, median capital of about 33,300 dollars and 880 watts
of facility power per H100-equivalent, with capital per H100-equivalent
ranging from about 19,600 to about 200,800 dollars {[}12{]}. These are
estimates built from public permits and imagery, and a site's listed
capital includes compute, construction and fit-out.


\begin{center}
\openiefig{fig/latency-fig09-small-multiples.png}{0.94}{https://research.openie.dev/figures/latency/small-multiples.png}
\end{center}


\emph{Figure 9. Four indicators in one panel. Capital and facility watts
per H100-equivalent for 74 sites sorted by value, peak tensor efficiency
of twelve accelerators in release order, and listed power in the six
largest countries. Data:
\texttt{latency\_\allowbreak{}data/\allowbreak{}frontier\_\allowbreak{}dc\_\allowbreak{}cost\_\allowbreak{}per\_\allowbreak{}h100e.\allowbreak{}csv},
\texttt{hardware\_\allowbreak{}flops\_\allowbreak{}per\_\allowbreak{}watt.\allowbreak{}csv},
\texttt{frontier\_\allowbreak{}dc\_\allowbreak{}by\_\allowbreak{}country.\allowbreak{}csv} {[}12, 13{]}.}

\textbf{Model output.} With lives of 4 to 6 years, utilization of 40\%
to 80\% and electricity at 0.04 to 0.12 dollars per kilowatt-hour, an
utilized accelerator-hour costs between about 0.83 and 2.5 dollars, and
energy is between 1.5\% and 11.8\% of it. The capital-loaded price
\(a/P\) is about 0.9 to 2.7 dollars per kilowatt-hour equivalent,
against electricity at 0.04 to 0.12, so \(\kappa^*\) is dominated by
capital at roughly 8 to 70 times the electricity term.


\begin{center}
\openiefig{fig/latency-fig10-accelerator-hour-cost.png}{0.94}{https://research.openie.dev/figures/latency/accelerator-hour-cost.png}
\end{center}


\emph{Figure 10. What one accelerator-hour costs and how little of it is
energy (model output). Capital 33,311 dollars and 880 W per
H100-equivalent, electricity 0.08 dollars per kilowatt-hour. Data:
\texttt{latency\_\allowbreak{}data/\allowbreak{}model\_\allowbreak{}outputs.\allowbreak{}json}, computed by
\texttt{model\_\allowbreak{}airace.\allowbreak{}py} from {[}12{]}.}

Three consequences follow. Cheaper electricity can lower the cost of a
task by at most the energy share, which is at most about an eighth in
this grid. Utilization matters as much as energy: halving \(u\) doubles
\(a\). And the stated life \(y\) is an accounting choice that moves
\(a\) in proportion to \(1/y\).

\textbf{Lives and efficiency.} Peak tensor operations per watt rose
about 1.42 times a year across twelve leading accelerators from 2016 to
2024 in the vendors' own specifications, which is a doubling about every
two years {[}13{]}. A power-limited site that replaced a part with one
three years newer would obtain about \(1.42^3 = 2.9\) times the peak
operations from the same power. Peak specifications are not delivered
energy per task, because utilization, memory traffic and software decide
the delivered figure. Stated server lives in Figure 12 run from 5 to 6
years at several large firms, and the range stated is wider at others.
Lives of that length and a two-year efficiency doubling are consistent
when older parts keep earning at prices set by newer ones, which is a
claim about prices and not about physics.


\begin{center}
\openiefig{fig/latency-fig11-hardware-efficiency.png}{0.94}{https://research.openie.dev/figures/latency/hardware-efficiency.png}
\end{center}


\emph{Figure 11. Peak tensor FLOP/s per watt of leading accelerators on
a log scale, with a fit of about 1.42 times a year. Vendor peak
specifications, not delivered energy per task. Data:
\texttt{latency\_\allowbreak{}data/\allowbreak{}hardware\_\allowbreak{}flops\_\allowbreak{}per\_\allowbreak{}watt.\allowbreak{}csv} and
\texttt{hardware\_\allowbreak{}fit.\allowbreak{}json} {[}13{]}.}


\begin{center}
\openiefig{fig/latency-fig12-depreciation-lives.png}{0.94}{https://research.openie.dev/figures/latency/depreciation-lives.png}
\end{center}


\emph{Figure 12. Stated useful lives for servers and network equipment,
in years, from annual reports. Lives are accounting estimates, not
measured service lives. Data: annual reports {[}4, 14, 15, 16, 17{]}.}

\textbf{Price is not cost.} For a system at the level of GPT-3.5, the
price of inference fell more than 280-fold between November 2022 and
October 2024, while hardware costs fell 30\% a year and hardware energy
efficiency improved 40\% a year {[}18{]}. Over the 1.9 years in that
window those rates give about 2.0 times for hardware cost and 1.9 times
for energy efficiency (model output). Dividing 280 by 2.0 leaves a
factor of about 141 that hardware cost does not explain. That remainder
mixes smaller models, algorithmic efficiency, utilization and margins,
and the sources do not apportion it. Epoch AI separately estimates that
the compute needed for a fixed language-model performance halves about
every 8 months {[}58{]}, and that the price for a fixed benchmark level
fell 9 to 900 times a year with a median near 50, the fastest rates
resting on under a year of data {[}59{]}.


\begin{center}
\openiefig{fig/latency-fig13-price-decomposition.png}{0.94}{https://research.openie.dev/figures/latency/price-decomposition.png}
\end{center}


\emph{Figure 13. A 280-fold price fall and how much hardware can explain
(model output from the annual rates in {[}18{]}). Data:
\texttt{latency\_\allowbreak{}data/\allowbreak{}model\_\allowbreak{}outputs.\allowbreak{}json}.}

\textbf{Tie to the frontiers.} Let \(\tau_C(t)\) be the lowest
accelerator time per task at quality \(q\) among documented methods and
\(\tau_D(t)\) the time in generally available products. Then
\(c_C(t) = \tau_C(t)\,\kappa^*(t)\) for the capability frontier, and the
deployed frontier follows the price offered, which is at least cost plus
the margin of Section 6. Hardware cost enters \(C_f\) and \(D\) in the
same way, because suppliers and buyers purchase from the same market, so
it shifts both frontiers together and leaves \(L\) nearly unchanged.
Latency is moved by the software term \(\tau\): a method that cuts
accelerator time per task by a factor \(g\) lowers \(c_C\) by that
factor when it is published, and \(c_D\) only when a product adopts it.
During the wait, the ratio of the deployed to the capability cost is
\(g\).

\textbf{Tie to the cadence model.} Capital is committed in multi-year
steps and charged over lives of several years, so the decision dates for
hardware are far apart, and the adoption cost \(K\) of Section 6
includes the value of installed capital that a new method would strand
or would save. A method that runs on hardware a firm already owns has an
adoption cost that is mostly engineering. A method that needs new
hardware carries the capital cycle with it. That is why open software
methods on existing parts are the least gated, and why Prediction 3
compares targets with and without porting tools.

\textbf{Where cost stops tracking capability.} At a fixed quality \(q\),
cost per task falls with \(\tau\), with hardware and with price. At the
frontier of capability, cost per task is set by the training and serving
of the largest models and does not fall in the same way, which is why a
frontier-quality task and a fixed-quality task have different cost
curves. The data above measure only the first.

\textbf{Prediction 11 (older parts keep earning).} If stated server
lives are economically realistic, rental prices per accelerator-hour for
parts several years old fall more slowly than peak operations per watt
would imply. \emph{Test:} dated open price listings by part and age,
against the efficiency fit in Figure 11. \emph{Refuted if} the price of
an older part per unit of peak operations falls at least as fast as the
efficiency fit.

\textbf{Prediction 12 (energy share stays small).} In a sample of
measured deployments the energy share of cost per task stays below 15\%
at electricity prices under 0.12 dollars per kilowatt-hour. \emph{Test:}
open energy measurement of inference {[}56{]} combined with listed
accelerator-hour prices. \emph{Refuted if} the measured share exceeds
15\% in a majority of deployments.

\section{13. Open methods and existing hardware at
parity}\label{open-methods-and-existing-hardware-at-parity}

A second description of the present is that open algorithmic work lets
existing and older infrastructure perform at parity with the newest. The
dated record supports a narrower statement, set out here by the lever
each method pulls and by the lag to a shipped use.

\textbf{Table 6. Open methods, their dated papers and what a later use
shows.}

{\def\LTcaptype{none} % do not increment counter
{\footnotesize\setlength{\tabcolsep}{3.8pt}\renewcommand{\arraystretch}{1.15}
\begin{longtable}[]{@{}
  >{\raggedright\arraybackslash}p{(\linewidth - 8\tabcolsep) * \real{0.1407}}
  >{\raggedright\arraybackslash}p{(\linewidth - 8\tabcolsep) * \real{0.3852}}
  >{\raggedright\arraybackslash}p{(\linewidth - 8\tabcolsep) * \real{0.0815}}
  >{\raggedright\arraybackslash}p{(\linewidth - 8\tabcolsep) * \real{0.2593}}
  >{\raggedright\arraybackslash}p{(\linewidth - 8\tabcolsep) * \real{0.1333}}@{}}
\toprule\noalign{}
\begin{minipage}[b]{\linewidth}\raggedright
\textbf{Lever}
\end{minipage} & \begin{minipage}[b]{\linewidth}\raggedright
\textbf{Open method and what the authors report}
\end{minipage} & \begin{minipage}[b]{\linewidth}\raggedright
\textbf{Posted}
\end{minipage} & \begin{minipage}[b]{\linewidth}\raggedright
\textbf{Later use}
\end{minipage} & \begin{minipage}[b]{\linewidth}\raggedright
\textbf{Lag}
\end{minipage} \\
\midrule\noalign{}
\endhead
\bottomrule\noalign{}
\endlastfoot
Quantization & Models with 175 billion parameters quantized in about
four GPU hours to 3 or 4 bits per weight with negligible accuracy loss
{[}92{]} & Oct 2022 & 4-bit local inference reported as practical on
consumer machines in 2023 {[}68{]} & Not computed: the later use is a
different 4-bit implementation, and its date is opinion-class
evidence \\
Ternary weights on CPU & A 2 billion parameter model at 0.4 GB of
non-embedding memory and 29 ms per token on a CPU, against 1.4 to 4.8 GB
and 41 to 124 ms for six open full-precision models of similar size
{[}26{]} & Apr 2025 & Model weights released with the paper & 0
months \\
Attention kernels & Exact attention made faster by accounting for reads
and writes between memory levels {[}19{]} & May 2022 & Accelerated
attention kernels in the PyTorch 2.0 release {[}20{]} & About 10
months \\
Serving memory & Paged key-value memory in an open serving system, with
higher throughput than prior systems at the same latency {[}93{]} & Sep
2023 & Source code released with the paper & 0 months \\
State-space and hybrid models & An attention-free model reporting 5
times the inference throughput of Transformers {[}21{]}; a hybrid that
fits one 80 GB GPU, with open weights {[}22{]} & Dec 2023 & The hybrid,
Mar 2024; more hybrids in 2024 and 2025 {[}39{]} & About 3 to 4
months \\
Speculative decoding & Output distribution unchanged with 2 to 3 times
faster decoding {[}23{]} & Nov 2022 & Reported in Google products by Dec
2024 {[}24{]} & About 24 months \\
Training cost & 2.788 million H800 GPU-hours for full training,
self-reported {[}46{]}; reasoning by reinforcement learning with
distilled smaller models {[}94{]} & Dec 2024 to Jan 2025 & Smaller
distilled models released with the second report & Not applicable \\
Non-NVIDIA serving & Prefill efficiency of 4.45 tokens per second per
TFLOPS on Ascend NPUs, above published results for SGLang on H100 and
DeepSeek on H800, power not reported {[}95{]} & Jun 2025 & Authors'
deployment & Not applicable \\
\end{longtable}
}
}

The lags are computed from the posting month to the date of the later
use, and rest on one pair per row. They are a sample of a few, not a
distribution, and the protocol of Section 16 specifies how to collect
many.


\begin{center}
\openiefig{fig/latency-fig14-research-to-product-lag.png}{0.94}{https://research.openie.dev/figures/latency/research-to-product-lag.png}
\end{center}


\emph{Figure 14. Paper to product, three dated pairs. First public paper
month and the date the same method is reported inside a shipped system.
Months are rounded and the sample is three. Data: dates in Table 6
{[}19, 20, 21, 22, 23, 24{]}.}

\textbf{Open models against the closed frontier, by training compute.}
Using the first notable model at each level of training compute, any
access type against open weights, the open-weight lag was 1 month at
\(10^{22}\) FLOP, 30 months at \(10^{23}\), 24 months at \(10^{24}\) and
15 months at \(10^{25}\), and no open-weight model is yet recorded at
\(10^{26}\) {[}25{]}. Each level is crossed first by one model, so the
lags rest on one pair per level. The direction is that the lag shrinks
as the level rises, and the caveat is that training compute is not
capability.


\begin{center}
\openiefig{fig/latency-fig15-open-compute-lag.png}{0.94}{https://research.openie.dev/figures/latency/open-compute-lag.png}
\end{center}


\emph{Figure 15. When open-weight models reached each training-compute
level. First notable model at or above the level, any access type
against open weights. Data:
\texttt{latency\_\allowbreak{}data/\allowbreak{}compute\_\allowbreak{}threshold\_\allowbreak{}firsts.\allowbreak{}csv}, from the Epoch AI
notable-models dataset {[}25{]}.}

On consumer hardware, open models that fit one consumer GPU match the
frontier after an average of 6 to 12 months, depending on the measure
{[}63{]}. Open-weight models of any size lag the closed frontier by
about one year on benchmarks, with a range of 5 to 22 months {[}64{]}.

\textbf{What parity means here.} It means the same benchmark score. It
does not mean the same throughput: a consumer card runs one stream, and
a data center serves many, where batching and kernels matter. The gains
from quantization and kernels apply on both. The compute needed for a
fixed performance halves about every 8 months {[}58{]}, which is
\(2^{12/8} = 2.8\) times a year, about twice the 1.42 times a year of
hardware efficiency in Figure 11. Where that holds, algorithmic progress
outruns hardware progress at a fixed quality, and parts that were bought
earlier produce more tasks of a given quality than they did. Whether
those extra tasks are sold at lower prices or are demanded in larger
numbers is a question about demand that none of the sources here
settles.


\begin{center}
\openiefig{fig/latency-fig16-timeline-capability-vs-deployment.png}{0.94}{https://research.openie.dev/figures/latency/timeline-capability-vs-deployment.png}
\end{center}


\emph{Figure 16. Capability releases against deployment and financing,
2022 to 2026. Numbered markers refer to the legend. Upper lane: public
methods and models. Lower lane: product releases, contracts and filings.
Data: dates of the sources in the References.}

\section{14. Other ways to build the
machine}\label{other-ways-to-build-the-machine}

A third description is that groups making hardware by orthogonal
processes show some of the cost of conventional infrastructure to be
higher than the task needs. The sources read contain dated
demonstrations of analog, photonic, neuromorphic, FPGA, CPU-only and
flexible-electronics designs. They differ in task, scale and
measurement, so this section reports each within its own paper and draws
only ratios that a paper itself reports.

\textbf{Table 7. Alternative designs, what was measured and the caveat.}

{\def\LTcaptype{none} % do not increment counter
{\footnotesize\setlength{\tabcolsep}{3.8pt}\renewcommand{\arraystretch}{1.15}
\begin{longtable}[]{@{}
  >{\raggedright\arraybackslash}p{(\linewidth - 6\tabcolsep) * \real{0.1544}}
  >{\raggedright\arraybackslash}p{(\linewidth - 6\tabcolsep) * \real{0.3020}}
  >{\raggedright\arraybackslash}p{(\linewidth - 6\tabcolsep) * \real{0.2819}}
  >{\raggedright\arraybackslash}p{(\linewidth - 6\tabcolsep) * \real{0.2617}}@{}}
\toprule\noalign{}
\begin{minipage}[b]{\linewidth}\raggedright
\textbf{Approach}
\end{minipage} & \begin{minipage}[b]{\linewidth}\raggedright
\textbf{Demonstration}
\end{minipage} & \begin{minipage}[b]{\linewidth}\raggedright
\textbf{Measured or reported}
\end{minipage} & \begin{minipage}[b]{\linewidth}\raggedright
\textbf{Caveat}
\end{minipage} \\
\midrule\noalign{}
\endhead
\bottomrule\noalign{}
\endlastfoot
Analog in-memory & 35 million phase-change memory devices across 34
tiles on a speech keyword task {[}96{]} & Up to 12.4 TOPS per watt
chip-sustained, accuracy equivalent to software & Small network, no
matched baseline in the text read, so no ratio is drawn \\
Photonic & A photonic AI processor at precision the authors present as
practical for ResNet and BERT {[}97{]} & Precision on those networks &
No joules per task in the text read \\
Neuromorphic & A 370 million parameter MatMul-free model on Loihi 2
{[}28{]} & Up to 3 times the throughput with 2 times less energy than
transformers on an edge GPU & Preliminary, small model \\
FPGA & FlightLLM on an Alveo U280 card running LLaMA2-7B with compiled
sparsity and quantization {[}27{]} & 6.0 and 4.2 times the energy
efficiency and 1.8 and 1.4 times the cost efficiency of a V100S and an
A100 & Authors' comparison and GPU configuration \\
Ternary weights on CPU & BitNet b1.58 2B on a CPU {[}26{]} & 0.4 GB
memory against 1.4 to 4.8 GB, and 29 ms against 41 to 124 ms per token &
The energy of 0.028 J per token is estimated by an arithmetic model, 6.6
to 23 times lower than the baselines, and was not read from a power
meter \\
Alternative accelerator stack & Serving on Ascend NPUs {[}95{]} & 4.45
tokens per second per TFLOPS in prefill, above published H100 and H800
results & Power not reported, so no joule ratio is drawn \\
Flexible, non-silicon process & A 32-bit Arm microprocessor on a
flexible substrate {[}98{]} & That a general processor can be made by a
non-silicon process & No machine-learning result and no joules per
task \\
\end{longtable}
}
}


\begin{center}
\openiefig{fig/latency-fig17-alternative-hardware-ratios.png}{0.94}{https://research.openie.dev/figures/latency/alternative-hardware-ratios.png}
\end{center}


\emph{Figure 17. Alternative designs against a baseline in the same
paper. Ratio is baseline cost divided by alternative cost, as each paper
reports it; tasks differ, so bars are not comparable across rows. Bars
show the range reported. Data: ratios in Table 7 {[}26, 27, 28{]}.}

\textbf{What this supports.} For particular task classes, namely small
models, edge deployment, CPU decoding and sparse or quantized inference
on programmable logic, the authors report energy or memory 2 to 23 times
lower than a conventional baseline. These are within-paper ratios on
small or mid-sized models. The sources read contain no measurement of an
alternative design serving a frontier-scale model, and none on
wafer-scale or mature-node designs, so this article does not draw a
conclusion for them.

\textbf{Why an energy ratio is not a cost ratio.} By Proposition 8, at
data-center prices energy is between about 1.5\% and 11.8\% of an
accelerator-hour. An alternative that lowers energy per task by a factor
of 10 lowers cost per task by at most the energy share, unless it also
lowers capital per task. The memory result of the ternary model points
at that second path: a model that needs 0.4 GB instead of 1.4 to 4.8 GB
can run on a smaller and cheaper device, and the capital effect has not
been measured in the sources here. Where electricity is not the dominant
term, the useful comparison is dollars per task including capital, and
none of the alternative-hardware papers read reports it except
FlightLLM, whose cost efficiency ratios are the smaller numbers in the
table.

\section{15. Where capability, compute and power
sit}\label{where-capability-compute-and-power-sit}

The geography of the build-out is read here from national statistics,
ministries and agencies, in the original language where the source is
not in English. Units differ across countries, so the figures are placed
side by side and are not ranked.

\textbf{Compute and power.} In the Epoch AI data-center dataset, listed
frontier sites total 14.1 GW, of which the United States holds 84\%,
followed by Malaysia at 1.1 GW and China at 0.7 GW {[}12{]}. The dataset
relies on public permits and imagery, so sites in places with less
public disclosure are more likely to be missed. China's own statistics
report a different quantity: 10.85 million standard racks in use and
intelligent computing capacity of 788 EFLOPS at the end of June 2025
{[}99{]}, and 1.59 million PFLOPS (FP16) at the end of 2025 {[}100{]}.
The two Chinese figures come from different publications and may rest on
different definitions, and neither is comparable with a listed-power
figure. For electricity, data centres used about 415 TWh in 2024, 1.5\%
of the world total, with 45\% in the United States, 25\% in China and
15\% in Europe, and the International Energy Agency projects a doubling
to about 945 TWh by 2030 as a modeled scenario {[}101, 102{]}.


\begin{center}
\openiefig{fig/latency-fig18-map-frontier-datacenters.png}{0.94}{https://research.openie.dev/figures/latency/map-frontier-datacenters.png}
\end{center}


\emph{Figure 18. Where frontier AI data-center power sits. Circle area
is proportional to listed power in megawatts; 94 sites, 14.1 GW. Data:
\texttt{latency\_\allowbreak{}data/\allowbreak{}frontier\_\allowbreak{}dc\_\allowbreak{}by\_\allowbreak{}country.\allowbreak{}csv} {[}12{]}.}

\textbf{Where models come from.} Of the notable models in 2025 by
national affiliation of the developer, the United States accounts for
59, China for 35 and South Korea for 8, with single models from Canada,
France, Hong Kong, Singapore and the United Kingdom {[}29{]}.


\begin{center}
\openiefig{fig/latency-fig19-map-notable-models.png}{0.94}{https://research.openie.dev/figures/latency/map-notable-models.png}
\end{center}


\emph{Figure 19. Where notable AI models came from in 2025, by national
affiliation of the developer. Data: Stanford AI Index 2026 as drawn from
the Epoch AI database {[}29{]}.}

\textbf{Chips.} For the accelerators in the Epoch hardware dataset, TSMC
is the recorded foundry for about four in five of the entries whose
foundry is recorded, followed by Samsung, SMIC and Intel, and the
foundry is not recorded for more than a third of all entries
(\texttt{latency\_\allowbreak{}data/\allowbreak{}hardware\_\allowbreak{}foundry\_\allowbreak{}counts.\allowbreak{}json}) {[}13{]}. The
sources read gave no usable figures for packaging capacity or for
fabrication capacity by country, so this article draws no map of them.

\textbf{National programs.} Programs state their size in different
units, and Table 8 gives each as stated.

\textbf{Table 8. National programs as stated, with the language of the
source.}

{\def\LTcaptype{none} % do not increment counter
{\footnotesize\setlength{\tabcolsep}{3.8pt}\renewcommand{\arraystretch}{1.15}
\begin{longtable}[]{@{}
  >{\raggedright\arraybackslash}p{(\linewidth - 6\tabcolsep) * \real{0.1667}}
  >{\raggedright\arraybackslash}p{(\linewidth - 6\tabcolsep) * \real{0.5625}}
  >{\raggedright\arraybackslash}p{(\linewidth - 6\tabcolsep) * \real{0.1146}}
  >{\raggedright\arraybackslash}p{(\linewidth - 6\tabcolsep) * \real{0.1562}}@{}}
\toprule\noalign{}
\begin{minipage}[b]{\linewidth}\raggedright
\textbf{Place}
\end{minipage} & \begin{minipage}[b]{\linewidth}\raggedright
\textbf{Statement}
\end{minipage} & \begin{minipage}[b]{\linewidth}\raggedright
\textbf{Language}
\end{minipage} & \begin{minipage}[b]{\linewidth}\raggedright
\textbf{Class}
\end{minipage} \\
\midrule\noalign{}
\endhead
\bottomrule\noalign{}
\endlastfoot
China & Intelligent computing capacity of 788 EFLOPS at end of June 2025
and 1.59 million PFLOPS (FP16) at end of 2025 {[}99, 100{]} & Chinese &
demonstrated \\
South Korea & Government statement of securing 260,000 NVIDIA GPUs, in
October 2025 {[}103{]} & Korean & opinion \\
Japan & At least 10 trillion yen of public support over the seven years
to fiscal 2030, aimed at more than 50 trillion yen of private investment
over ten years {[}104{]}; market research stating that cloud investment
in Japanese data centers is expected to be hundreds of billions of yen
or more each year {[}105{]} & Japanese & opinion \\
European Union & 200 billion euro to be mobilised for AI, with a 20
billion euro fund for AI gigafactories {[}106, 107{]} & English and
Italian & opinion \\
India & 38,231 GPUs onboarded under the IndiaAI compute portal, as
reported to parliament in March 2026 {[}108{]} & English &
demonstrated \\
Brazil & R\$ 23.03 billion of planned investment for 2024 to 2028
{[}109{]} \textbar{} Portuguese \textbar{} opinion & & \\
United Arab Emirates & A 1 GW compute cluster inside a planned 5 GW
campus in Abu Dhabi {[}110{]} & English & opinion \\
Saudi Arabia & Projected capacity of up to 500 MW over five years
{[}111{]} & English & opinion \\
Africa & An AI factory deploying in South Africa, with plans for
Nigeria, Kenya, Egypt and Morocco; no capacity figure stated {[}112{]} &
English & opinion \\
\end{longtable}
}
}


\begin{center}
\openiefig{fig/latency-fig20-map-national-programs.png}{0.94}{https://research.openie.dev/figures/latency/map-national-programs.png}
\end{center}


\emph{Figure 20. National compute programs and their stated sizes. Units
differ by country, so no ranking is implied. Data: Table 8.}

Statements in the Opinion class are plans and targets, and the
Demonstrated class holds reported capacity or statistics. The two kinds
should not be added.

\textbf{Use.} Official surveys of firms using AI give 32.9\% in South
Korea in 2024, 26\% in Germany in 2025, 21.1\% in Spain in early 2025,
20.0\% in the European Union in 2025, 18\% in the United States in the
winter of 2025 to 2026, 13.9\% for the OECD area in 2024, 13\% in Brazil
in 2024 and 10\% in France in 2024, with years and definitions that
differ {[}30, 31, 32, 33, 34, 35, 36, 37{]}. Japan and China are
measured on people and not on firms {[}71, 72{]}.


\begin{center}
\openiefig{fig/latency-fig21-adoption-by-country.png}{0.94}{https://research.openie.dev/figures/latency/adoption-by-country.png}
\end{center}


\emph{Figure 21. Share of firms using AI in official surveys. Years and
definitions differ, so each row is read on its own. Data: the
statistical releases cited in Section 3 {[}30, 31, 32, 33, 34, 35, 36,
37{]}.}

\textbf{What the geography shows.} Listed compute power is concentrated
in one country by a wide margin, while published models and official
capacity statements show a second country with a large share of models
and a large stated capacity that listed-power datasets are poorly placed
to see. Programs elsewhere state plans of gigawatt or hundred-billion
scale, and the measured use of AI by firms is between 10\% and 33\% in
the surveys read. The gap between plans and measured use is the gap this
article studies, now at national scale.

\section{16. Experiment protocol (not
executed)}\label{experiment-protocol-not-executed}

\textbf{Purpose.} Measure \(L\) and its three parts for a list of
methods, using only public dates.

\textbf{Unit.} A method, such as a number format, a kernel, an inference
procedure or a compiler pass, with a dated public description.

\textbf{Dates to collect for each method.} The earliest public posting
(preprint or repository), the earliest open implementation, the earliest
release in a widely used open library, and the earliest
general-availability release in at least one commercial product, with
the firm type (listed, private, public-sector, academic). Each date is
recorded with its source link and an archive copy.

\textbf{Sampling.} All methods in a pre-registered list drawn from a
fixed set of venues and repositories over 2022 to 2026, including
non-English venues and repositories, with the selection rule written
before any release date is looked up.

\textbf{Analysis.} Compute \(L\) per method. Estimate \(E[w^{cad}]\) by
the phase tests of Prediction 1 and by comparing the mean delay with
\(T/2\). Estimate the gate part as the difference between listed and
unlisted firms (Prediction 2). Report medians with bootstrap intervals,
and report what share of methods never reach a product.

\textbf{Energy.} Where a method reports joules per task, record the
measurement conditions, and where it does not, run it on an open energy
harness such as the one cited in Table 1 and report the measured joules.

\textbf{Financing and cost.} Compute the bound of Definition 4 for each
supplier and fiscal year from filings, the capital-loaded price of a
joule from listed accelerator-hour prices and measured power, and the
ratio of contracted purchases to trailing revenue for each lab, and
record each with the filing and page it came from. The figure script
\texttt{scripts/\allowbreak{}figures/\allowbreak{}latency\_\allowbreak{}figures.\allowbreak{}py} and the data files under
\texttt{scripts/\allowbreak{}figures/\allowbreak{}latency\_\allowbreak{}data/\allowbreak{}} reproduce every figure in
Sections 9 to 15 from the cited sources.

\textbf{Outputs.} One table of dates, one analysis script, and a log of
exclusions. The scripts and data are to be released with the study under
the same terms as this site.

\section{17. Conclusion}\label{conclusion}

Public methods that lower the cost of computer intelligence reach
products after a delay with a technical part, a calendar part and an
incentive part. The first is engineering time. The second is the
structure of decision dates, which on a quarterly calendar adds about
six weeks on average (Proposition 1). The third appears when a reporting
rule blocks a release that is worth doing in the long run (Proposition
3). A tool that lowers the cost of adopting a method removes the third
once the cost falls under a stated bound, and a tool that lowers the
cost of the work itself forces release when holding earns nothing
(Proposition 4). When the market has several sellers reachable at low
friction, the fastest sets the latency (Proposition 5). As computing
cost falls, the margin that remains depends on switching friction
(Proposition 6). The infrastructure evidence adds the following. Filings
document loops of tens of billions of dollars between named suppliers
and customers, and they bound the share of the largest chip supplier's
revenue that its own capital could have funded under a fifth, while most
hyperscaler spending is paid from operating cash flow. At data-center
prices energy is a small part of an accelerator-hour and capital is the
large part, so cost per task follows accelerator time per task,
utilization and the stated life of the machine. Open methods and
alternative designs reduce accelerator time, memory or energy per task
in the authors' own measurements, for small and mid-sized models, and no
source read measures them at frontier scale. Capability, compute and
power are located unevenly, and national plans are stated at scales well
above measured use. The propositions are a model, and the predictions
name the filings, listings and measurements that would refute each.

\section{References}\label{references}

Each entry ends with its evidence class (Section 1).

\begin{enumerate}
\def\labelenumi{\arabic{enumi}.}
\tightlist
\item
  Meta, joint venture with funds managed by Blue Owl Capital to develop
  the Hyperion data center. 2025-10-21.
  \url{https://investor.atmeta.com/investor-news/press-release-details/2025/Meta-Announces-Joint-Venture-with-Funds-Managed-by-Blue-Owl-Capital-to-Develop-Hyperion-Data-Center/default.aspx}
  (opinion)
\item
  Meta Platforms, Quarterly Report on Form 10-Q, period ended 30
  September 2025. 2025-10-30.
  \url{https://www.sec.gov/Archives/edgar/data/1326801/000162828025047240/meta-20250930.htm}
  (demonstrated)
\item
  NVIDIA, Annual Report on Form 10-K, fiscal year ended 25 January 2026.
  2026-02-25.
  \url{https://www.sec.gov/Archives/edgar/data/1045810/000104581026000021/nvda-20260125.htm}
  (demonstrated)
\item
  CoreWeave, Annual Report on Form 10-K, fiscal year ended 31 December
  2025. 2026-03-02.
  \url{https://www.sec.gov/Archives/edgar/data/1769628/000176962826000104/crwv-20251231.htm}
  (demonstrated)
\item
  Oracle, Quarterly Report on Form 10-Q, period ended 31 August 2026.
  2026-09-11.
  \url{https://www.sec.gov/Archives/edgar/data/1341439/000119312626389274/orcl-20260831.htm}
  (demonstrated)
\item
  Oracle, Q1 FY27 results (Exhibit 99.1 to Form 8-K). 2026-09-10.
  \url{https://www.sec.gov/Archives/edgar/data/1341439/000119312626387905/orcl-20260910.htm}
  (demonstrated)
\item
  Microsoft, Quarterly Report on Form 10-Q, period ended 30 September
  2025, Note 17. 2025-10-29.
  \url{https://www.sec.gov/Archives/edgar/data/789019/000119312525256321/msft-20250930.htm}
  (demonstrated)
\item
  OpenAI and Amazon Web Services, partnership announcement. 2025-11-03.
  \url{https://openai.com/index/aws-and-openai-partnership/} (opinion)
\item
  CNBC, OpenAI closes record-breaking 122 billion dollar funding round.
  2026-03-31.
  \url{https://www.cnbc.com/2026/03/31/openai-funding-round-ipo.html}
  (opinion)
\item
  Microsoft, NVIDIA and Anthropic, strategic partnerships announcement.
  2025-11-18.
  \url{https://blogs.microsoft.com/blog/2025/11/18/microsoft-nvidia-and-anthropic-announce-strategic-partnerships/}
  (opinion)
\item
  U.S. Securities and Exchange Commission, XBRL company concept API
  (annual capital expenditure and operating cash flow, seven filers).
  2026-10. \url{https://data.sec.gov/api/xbrl/companyconcept/}
  (demonstrated)
\item
  Epoch AI, Frontier Data Centers dataset (CC BY). 2026-10.
  \url{https://epoch.ai/data/data-centers} (demonstrated)
\item
  Epoch AI, Machine Learning Hardware dataset (CC BY). 2026-10.
  \url{https://epoch.ai/data/machine-learning-hardware} (demonstrated)
\item
  Amazon, Annual Report on Form 10-K, fiscal year ended 31 December
  2025. 2026-02-06.
  \url{https://www.sec.gov/Archives/edgar/data/1018724/000101872426000004/amzn-20251231.htm}
  (demonstrated)
\item
  Meta Platforms, Annual Report on Form 10-K, fiscal year ended 31
  December 2025. 2026-01-29.
  \url{https://www.sec.gov/Archives/edgar/data/1326801/000162828026003942/meta-20251231.htm}
  (demonstrated)
\item
  Alphabet, Annual Report on Form 10-K, fiscal year ended 31 December
  2025. 2026-02-05.
  \url{https://www.sec.gov/Archives/edgar/data/1652044/000165204426000018/goog-20251231.htm}
  (demonstrated)
\item
  Microsoft, Annual Report on Form 10-K, fiscal year ended 30 June 2026.
  2026-07-29.
  \url{https://www.sec.gov/Archives/edgar/data/789019/000119312626323660/msft-20260630.htm}
  (demonstrated)
\item
  Stanford HAI, AI Index Report 2025. 2025-04.
  \url{https://hai.stanford.edu/assets/files/hai_ai_index_report_2025.pdf}
  (demonstrated)
\item
  Dao et al., FlashAttention: Fast and Memory-Efficient Exact Attention
  with IO-Awareness. 2022-05-27. \url{https://arxiv.org/abs/2205.14135}
  (demonstrated)
\item
  PyTorch, PyTorch 2.0 release announcement. 2023-03-15.
  \url{https://pytorch.org/blog/pytorch-2-0-release/} (opinion)
\item
  Gu and Dao, Mamba: Linear-Time Sequence Modeling with Selective State
  Spaces. 2023-12-01. \url{https://arxiv.org/abs/2312.00752}
  (demonstrated)
\item
  Lieber et al., Jamba: A Hybrid Transformer-Mamba Language Model.
  2024-03-28. \url{https://arxiv.org/abs/2403.19887} (demonstrated)
\item
  Leviathan, Kalman, Matias, Fast Inference from Transformers via
  Speculative Decoding. 2022-11-30 (ICML 2023).
  \url{https://arxiv.org/abs/2211.17192} (established by proof in the
  source)
\item
  Google Research, Looking back at speculative decoding. 2024-12-06.
  \url{https://research.google/blog/looking-back-at-speculative-decoding/}
  (opinion)
\item
  Epoch AI, Notable AI Models dataset (CC BY). 2026-10.
  \url{https://epoch.ai/data/notable-ai-models} (demonstrated)
\item
  Ma et al., BitNet b1.58 2B4T Technical Report. 2025-04-16.
  \url{https://arxiv.org/abs/2504.12285} (demonstrated)
\item
  Zeng et al., FlightLLM: Efficient Large Language Model Inference with
  a Complete Mapping Flow on FPGAs. 2024-01-08.
  \url{https://arxiv.org/abs/2401.03868} (demonstrated)
\item
  Abreu et al., Neuromorphic Principles for Efficient Large Language
  Models on Intel Loihi 2. 2025-03.
  \url{https://arxiv.org/abs/2503.18002} (demonstrated)
\item
  Stanford HAI, AI Index Report 2026, research and development chapter.
  2026-04. \url{https://hai.stanford.edu/ai-index/2026-ai-index-report}
  (demonstrated)
\item
  (Korean) Electronic Times, reporting the 2025 Enterprise
  Informatization Statistics of the Ministry of Science and ICT and NIA.
  2026-01-20. \url{https://www.etnews.com/20260120000071} (demonstrated)
\item
  (German) Statistisches Bundesamt, Unternehmen mit Nutzung von
  Technologien der künstlichen Intelligenz nach
  Beschäftigtengrößenklassen (table, 2025). 2025.
  \url{https://www.destatis.de/DE/Themen/Branchen-Unternehmen/Unternehmen/IKT-in-Unternehmen-IKT-Branche/Tabellen/ikti-unternehmen-kuenstliche-intelligenz.html}
  (demonstrated)
\item
  (Spanish) Instituto Nacional de Estadística, Encuesta sobre el uso de
  TIC y del comercio electrónico en las empresas, nota de prensa.
  2025-10-22. \url{https://ine.es/dyngs/Prensa/ETICCE20241T2025.pdf}
  (demonstrated)
\item
  Eurostat, 20\% of EU enterprises use AI technologies. 2025-12-11.
  \url{https://ec.europa.eu/eurostat/documents/7870049/23260410/KS-01-26-009-EN-N.pdf}
  (demonstrated)
\item
  Bonney et al., The Microstructure of AI Diffusion (Census
  CES-WP-26-25). 2026-04.
  \url{https://www.census.gov/library/working-papers/2026/adrm/CES-WP-26-25.html}
  (demonstrated)
\item
  OECD, Emerging divides in the transition to artificial intelligence.
  2025-06.
  \url{https://www.oecd.org/content/dam/oecd/en/publications/reports/2025/06/emerging-divides-in-the-transition-to-artificial-intelligence_eeb5e120/7376c776-en.pdf}
  (demonstrated)
\item
  (Portuguese) Cetic.br/NIC.br, Pesquisa TIC Empresas 2024, resumo
  executivo. 2025-05.
  \url{https://www.nic.br/media/docs/publicacoes/2/20250512121759/tic_empresas_2024_resumo_executivo.pdf}
  (demonstrated)
\item
  (French) Insee, Les technologies de l'information et de la
  communication dans les entreprises en 2024, Insee Première n° 2061.
  2025-07.
  \url{https://www.insee.fr/fr/statistiques/fichier/8604126/ip2061.pdf}
  (demonstrated)
\item
  Dao and Gu, Transformers are SSMs (Mamba-2, state space duality).
  2024-05-31. \url{https://arxiv.org/abs/2405.21060} (established by
  proof in the source)
\item
  Lieber et al.~(Jamba); NVIDIA (Nemotron-H); TII (Falcon-H1); IBM
  (Granite 4.0). 2024 to 2025. \url{https://arxiv.org/abs/2504.03624}
  (demonstrated)
\item
  Ma et al., The Era of 1-bit LLMs (BitNet b1.58). 2024-02-27.
  \url{https://arxiv.org/abs/2402.17764} (demonstrated)
\item
  NVIDIA, Pretraining Large Language Models with NVFP4. 2025-09.
  \url{https://arxiv.org/abs/2509.25149} (demonstrated)
\item
  Shah et al., FlashAttention-3. 2024-07.
  \url{https://arxiv.org/abs/2407.08608} (demonstrated)
\item
  Yuan et al., Native Sparse Attention. 2025-02 (ACL 2025).
  \url{https://arxiv.org/abs/2502.11089} (demonstrated)
\item
  Hoffmann et al., Training Compute-Optimal Large Language Models.
  2022-03-29. \url{https://arxiv.org/abs/2203.15556} (demonstrated)
\item
  Sardana et al., Beyond Chinchilla-Optimal: Accounting for Inference in
  Language Model Scaling. 2024-01 (ICML 2024).
  \url{https://arxiv.org/abs/2401.00448} (demonstrated)
\item
  DeepSeek-AI, DeepSeek-V3 Technical Report. 2024-12-26.
  \url{https://arxiv.org/abs/2412.19437} (demonstrated)
\item
  Angelopoulos, Barber, Bates, Theoretical Foundations of Conformal
  Prediction. 2024-11-18 (v5 2026-03-06).
  \url{https://arxiv.org/abs/2411.11824} (established by proof in the
  source)
\item
  Zhou, Chen, Gui and Cheng, Conformal Prediction: A Data Perspective.
  2024-10 (arXiv 2410.06494; ACM Computing Surveys 2025).
  \url{https://arxiv.org/abs/2410.06494} (opinion)
\item
  Angelopoulos et al., Prediction-powered inference. 2023 (Science).
  \url{https://www.science.org/doi/10.1126/science.adi6000} (established
  by proof in the source)
\item
  Ramdas, Grunwald, Vovk, Shafer, Game-Theoretic Statistics and Safe
  Anytime-Valid Inference. 2023-11 (Statistical Science 38(4)).
  \url{https://doi.org/10.1214/23-sts894} (established by proof in the
  source)
\item
  Chernozhukov et al., Double/debiased machine learning. 2018
  (Econometrics Journal). \url{https://doi.org/10.1111/ectj.12097}
  (established by proof in the source)
\item
  Jung, Tian, Bareinboim, Estimating Identifiable Causal Effects through
  Double Machine Learning. 2021 (AAAI).
  \url{https://doi.org/10.1609/aaai.v35i13.17438} (established by proof
  in the source)
\item
  Sathia et al., Compiler Optimizations for Higher-Order Automatic
  Differentiation. 2025 (PPoPP workshop).
  \url{https://ece.is/assets/pdf/diffprog-ppopp25-paper90.pdf}
  (demonstrated)
\item
  Ward et al., JAX a la Stan. 2026-10-11.
  \url{https://bob-carpenter.github.io/dj-paper/template-computo-python.pdf}
  (opinion)
\item
  Qian et al., Impulse: probabilistic programming via compiler
  augmentation. 2026 (abstract).
  \url{https://ece.is/assets/pdf/impulse-pldi-src-abstract.pdf}
  (demonstrated)
\item
  Chung et al., The ML.ENERGY Benchmark. 2025 (NeurIPS Datasets and
  Benchmarks). \url{https://ml.energy/leaderboard/} (demonstrated)
\item
  Google, Measuring the environmental impact of delivering AI at Google
  Scale. 2025-08. \url{https://arxiv.org/abs/2508.15734} (demonstrated)
\item
  Ho et al.~(Epoch AI), Algorithmic progress in language models. 2024-03
  (NeurIPS 2024). \url{https://arxiv.org/abs/2403.05812} (demonstrated)
\item
  Epoch AI, LLM inference prices have fallen rapidly but unequally
  across tasks. 2025.
  \url{https://epoch.ai/data-insights/llm-inference-price-trends}
  (demonstrated)
\item
  Comin and Hobijn, An Exploration of Technology Diffusion. 2010 (AER
  100(5)). \url{https://doi.org/10.1257/aer.100.5.2031} (demonstrated)
\item
  P. A. David, The Dynamo and the Computer. 1990 (AER 80(2)).
  \url{https://ideas.repec.org/a/aea/aecrev/v80y1990i2p355-61.html}
  (demonstrated)
\item
  Brynjolfsson, Rock, Syverson, The Productivity J-Curve. 2021 (AEJ
  Macro). \url{https://doi.org/10.1257/mac.20180386} (demonstrated)
\item
  Somala and Emberson (Epoch AI), Frontier AI performance becomes
  accessible on consumer hardware within a year. 2025-08-15.
  \url{https://epoch.ai/data-insights/consumer-gpu-model-gap}
  (demonstrated)
\item
  Epoch AI, Open vs.~closed AI: how behind are open models?. 2025.
  \url{https://epoch.ai/publications/open-models-report} (demonstrated)
\item
  Gupta and Kiely, Production-ready speculative decoding with
  TensorRT-LLM. 2024-12-19.
  \url{https://www.baseten.co/blog/how-we-built-production-ready-speculative-decoding-with-tensorrt-llm/}
  (opinion)
\item
  Marx and Fuegi, Reliance on Science: Worldwide Front-Page Patent
  Citations to Scientific Articles. 2020 (NBER 27987; JEMS 2022).
  \url{https://www.nber.org/system/files/working_papers/w27987/w27987.pdf}
  (demonstrated)
\item
  Ahmadpoor and Jones, The dual frontier. 2017 (Science).
  \url{https://doi.org/10.1126/science.aam9527} (demonstrated)
\item
  Willison, Stanford Alpaca, and the acceleration of on-device large
  language model development. 2023-03-13.
  \url{https://simonwillison.net/2023/Mar/13/alpaca/} (opinion)
\item
  Li et al., Now Hiring: Firm-Level Adoption of Generative AI in China
  (World Bank). 2025/2026.
  \url{https://thedocs.worldbank.org/en/doc/cbdef4ef57efdbb3c29aa4c07edf860b-0050062026/original/Now-hiring-Firm-Level-Li-etal.pdf}
  (demonstrated)
\item
  (German) Statistisches Bundesamt, Pressemitteilung Nr. 444: Jedes
  fünfte Unternehmen nutzt künstliche Intelligenz. 2024-11-25.
  \url{https://www.destatis.de/DE/Presse/Pressemitteilungen/2024/11/PD24_444_52911.html}
  (demonstrated)
\item
  Japan Ministry of Internal Affairs and Communications,
  令和7年版情報通信白書 (Information and Communications White Paper
  2025), data tables, in Japanese. 2025.
  \url{https://www.soumu.go.jp/johotsusintokei/whitepaper/ja/r07/html/nd112220.html}
  (demonstrated)
\item
  (Chinese) China Internet Network Information Center,
  生成式人工智能应用发展报告（2025）. 2025-10.
  \url{https://pdf.dfcfw.com/pdf/H301_AP202510241768289458_1.pdf}
  (demonstrated)
\item
  Bick, Blandin, Deming, The Rapid Adoption of Generative AI. 2024-09
  (rev. 2025-02). \url{https://www.nber.org/papers/w32966}
  (demonstrated)
\item
  Brynjolfsson, Li, Raymond, Generative AI at Work (NBER Working Paper
  31161; later published in QJE 2025). 2023 working paper, revised.
  \url{https://www.nber.org/papers/w31161} (demonstrated)
\item
  METR, Measuring the Impact of Early-2025 AI on Experienced Open-Source
  Developer Productivity. 2025-07-10.
  \url{https://arxiv.org/abs/2507.09089} (demonstrated)
\item
  Bloom, Jones, Van Reenen, Webb, Are Ideas Getting Harder to Find?.
  2020 (AER 110(4)). \url{https://doi.org/10.1257/aer.20180338}
  (demonstrated)
\item
  Graham, Harvey, Rajgopal, The Economic Implications of Corporate
  Financial Reporting. 2005 (JAE 40); NBER 2004-10.
  \url{https://www.nber.org/papers/w10550} (demonstrated)
\item
  Asker, Farre-Mensa, Ljungqvist, Corporate Investment and Stock Market
  Listing: A Puzzle?. 2015 (RFS 28(2)).
  \url{https://doi.org/10.1093/rfs/hhu077} (demonstrated)
\item
  FactSet, Highest number of S\&P 500 earnings calls citing AI over the
  past 10 years. 2025-12.
  \url{https://insight.factset.com/highest-number-of-sp-500-earnings-calls-citing-ai-over-the-past-10-years-1}
  (demonstrated)
\item
  Hassan, Hollander, van Lent, Tahoun, Firm-Level Political Risk. 2019
  (QJE 134(4)). \url{https://doi.org/10.1093/qje/qjz021} (demonstrated)
\item
  Metrick and Yasuda, Venture Capital and Other Private Equity: A Survey
  (NBER Working Paper 16652). 2010-10.
  \url{https://www.nber.org/system/files/working_papers/w16652/w16652.pdf}
  (opinion)
\item
  Gompers and Lerner, What Drives Venture Capital Fundraising?
  (Brookings Papers: Microeconomics). 1998.
  \url{https://www.brookings.edu/wp-content/uploads/1998/01/1998_bpeamicro_gompers.pdf}
  (demonstrated)
\item
  Arrow, Economic Welfare and the Allocation of Resources for Invention
  (NBER chapter). 1962.
  \url{https://www.nber.org/system/files/chapters/c2144/c2144.pdf}
  (established by proof in the source)
\item
  Shapiro, Competition and Innovation: Did Arrow Hit the Bull's Eye?
  (NBER conference draft). 2010.
  \url{https://conference.nber.org/confer/2010/RDIA10/Shapiro.pdf}
  (opinion)
\item
  Henderson and Clark, Architectural Innovation. 1990 (ASQ 35(1)).
  \url{https://doi.org/10.2307/2393549} (demonstrated)
\item
  US District Court for D.C., United States v. Google LLC, remedies
  memorandum. 2025-09-02.
  \url{https://www.courthousenews.com/wp-content/uploads/2025/09/judge-amit-mehta-google-search-monopoly-remedy-ruling.pdf}
  (demonstrated)
\item
  NVIDIA and OpenAI, letter of intent announcement. 2025-09-22.
  \url{https://nvidianews.nvidia.com/news/openai-and-nvidia-announce-strategic-partnership-to-deploy-10gw-of-nvidia-systems}
  (opinion)
\item
  Advanced Micro Devices, Form 8-K, Item 1.01 (warrant issued to
  OpenAI). 2025-10-06.
  \url{https://www.sec.gov/Archives/edgar/data/2488/000119312525230895/d28189d8k.htm}
  (demonstrated)
\item
  Microsoft, The next chapter of the Microsoft-OpenAI partnership.
  2025-10-28.
  \url{https://blogs.microsoft.com/blog/2025/10/28/the-next-chapter-of-the-microsoft-openai-partnership/}
  (opinion)
\item
  Anthropic, Expanding our use of Google Cloud TPUs and services.
  2025-10-23.
  \url{https://www.anthropic.com/news/expanding-our-use-of-google-cloud-tpus-and-services}
  (opinion)
\item
  OpenAI and Broadcom, strategic collaboration announcement. 2025-10-13.
  \url{https://openai.com/index/openai-and-broadcom-announce-strategic-collaboration/}
  (opinion)
\item
  Frantar et al., GPTQ: Accurate Post-Training Quantization for
  Generative Pre-trained Transformers. 2022-10-31.
  \url{https://arxiv.org/abs/2210.17323} (demonstrated)
\item
  Kwon et al., Efficient Memory Management for Large Language Model
  Serving with PagedAttention (vLLM). 2023-09-12.
  \url{https://arxiv.org/abs/2309.06180} (demonstrated)
\item
  DeepSeek-AI, DeepSeek-R1: Incentivizing Reasoning Capability in LLMs
  via Reinforcement Learning. 2025-01-22.
  \url{https://arxiv.org/abs/2501.12948} (demonstrated)
\item
  Zuo et al., Serving Large Language Models on Huawei CloudMatrix384.
  2025-06. \url{https://arxiv.org/abs/2506.12708} (demonstrated)
\item
  Le Gallo et al., An analog-AI chip for energy-efficient speech
  recognition and transcription, Nature. 2023-08-23.
  \url{https://www.nature.com/articles/s41586-023-06337-5}
  (demonstrated)
\item
  Ahmed et al., Universal photonic artificial intelligence acceleration,
  Nature. 2025-04-09.
  \url{https://www.nature.com/articles/s41586-025-08854-x}
  (demonstrated)
\item
  Biggs et al., A natively flexible 32-bit Arm microprocessor, Nature.
  2021-07-21. \url{https://www.nature.com/articles/s41586-021-03625-w}
  (demonstrated)
\item
  (Chinese) 国务院, 截至今年6月底我国智能算力规模达788百亿亿次/秒
  (人民日报). 2025-09-02.
  \url{https://www.gov.cn/yaowen/liebiao/202509/content_7038732.htm}
  (demonstrated)
\item
  (Chinese) 全国数据资源统计调查工作组, 全国数据资源调查报告（2025年）.
  2026.
  \url{https://www.nda.gov.cn/sjj/swdt/xwfb/0429/ff808081-9b5e8626-019d-d86b2f54-1154.pdf}
  (demonstrated)
\item
  International Energy Agency, Energy and AI, executive summary.
  2025-04.
  \url{https://www.iea.org/reports/energy-and-ai/executive-summary}
  (demonstrated)
\item
  International Energy Agency, Energy and AI, energy demand from AI.
  2025-04.
  \url{https://www.iea.org/reports/energy-and-ai/energy-demand-from-ai}
  (demonstrated)
\item
  (Korean) 뉴시스, 정부, 엔비디아 GPU 26만장 확보. 2025-10-31.
  \url{https://www.newsis.com/view/NISX20251031_0003385059} (opinion)
\item
  (Japanese) 経済産業省, AI・半導体産業基盤強化フレーム. 2024-11.
  \url{https://www.meti.go.jp/policy/mono_info_service/ai_semiconductor_frame/ai_semiconductor_frame.html}
  (opinion)
\item
  (Japanese) 総務省 有識者会合 資料5,
  国内データセンター市場におけるAI需要・地方分散・再エネ電源
  (富士キメラ総研). 2024-05-30.
  \url{https://www.meti.go.jp/policy/mono_info_service/joho/conference/digital_infrastructure/0007/005_fujichimera.pdf}
  (opinion)
\item
  European Commission, EU launches InvestAI initiative to mobilise 200
  billion euro (English). 2025-02-11.
  \url{https://digital-strategy.ec.europa.eu/en/news/eu-launches-investai-initiative-mobilise-eu200-billion-investment-artificial-intelligence}
  (opinion)
\item
  (Italian) Commissione europea, UE lancia InvestAI per mobilitare 200
  miliardi di EUR. 2025-02-11.
  \url{https://digital-strategy.ec.europa.eu/it/news/eu-launches-investai-initiative-mobilise-eu200-billion-investment-artificial-intelligence}
  (opinion)
\item
  Press Information Bureau (India), reply in the Lok Sabha on the
  IndiaAI Mission. 2026-03-25.
  \url{https://www.pib.gov.in/PressReleasePage.aspx?PRID=2245063}
  (demonstrated)
\item
  (Portuguese) Ministério da Ciência, Tecnologia e Inovação, Plano
  Brasileiro de Inteligência Artificial 2024-2028. 2025.
  \url{https://www.gov.br/mcti/pt-br/acompanhe-o-mcti/transformacaodigital/plano-brasileiro-de-inteligencia-artificial}
  (opinion)
\item
  SoftBank Group, Global Tech Alliance Launches Stargate UAE.
  2025-05-22. \url{https://group.softbank/en/news/press/20250522}
  (opinion)
\item
  NVIDIA, HUMAIN and NVIDIA Announce Strategic Partnership to Build AI
  Factories in Saudi Arabia. 2025-05-13.
  \url{https://nvidianews.nvidia.com/news/humain-and-nvidia-announce-strategic-partnership-to-build-ai-factories-of-the-future-in-saudi-arabia}
  (opinion)
\item
  Cassava Technologies, Cassava scales African AI infrastructure with
  NVIDIA-powered AI factories. 2025.
  \url{https://www.cassavatechnologies.com/cassava-scales-african-ai-infrastructure-with-nvidia-powered-ai-factories-to-accelerate-sovereign-data-capabilities/}
  (opinion)
\end{enumerate}


\end{document}
